\section{Part A: cohomological dimension}
\label{A:section}

\subsection{Conventions and two basic comparisons}

Throughout this part, $\cd_\ell(T)$ refers to cohomology of \emph{all} étale $\ell$-primary torsion abelian sheaves on $T$. In particular, testing constant sheaves alone does not compute this invariant. An infinite right-hand side makes an upper bound vacuous. A field of characteristic $\ell$ has étale $\ell$-cohomological dimension at most one; its $\ell$-rank and its Kato $\ell$-dimension are different invariants. We write $d_\ell(F)$ for Kato's dimension in characteristic $\ell$, and for $\cd_\ell(F)$ in other characteristics. For a characteristic-$\ell$ field $F$, we write $\delta(F)$ for the cardinality of an absolute $\ell$-basis, interpreting any infinite value as $+\infty$ in numerical inequalities.

\begin{lemma}[Known; coefficient and closed-immersion comparisons]
\label{A:basic}
Let $T$ be noetherian. Uniform vanishing for constructible $\mathbf F_\ell$-sheaves in degrees above $N$ implies $\cd_\ell(T)\leq N$. If $i:Z\hookrightarrow T$ is a closed immersion, then $\cd_\ell(Z)\leq\cd_\ell(T)$.
\end{lemma}
\begin{proof}
On a noetherian scheme every torsion sheaf is a filtered union of constructible subsheaves, and étale cohomology commutes with the corresponding filtered colimits. The usual dévissage by multiplication by $\ell$ reduces bounded-exponent coefficients to $\mathbf F_\ell$-coefficients; taking the union over the exponents gives the stated assertion. These are the constructible-approximation and continuity statements of \citet[Exposés VI, \S5, and IX, \S2]{A_SGA4}. For the second assertion, $i_*$ on étale abelian sheaves is exact, its higher direct images vanish, and $H^q(T,i_*F)=H^q(Z,F)$ for every $q$ and every étale abelian sheaf $F$ on $Z$.
\end{proof}

\begin{proposition}[Proved; generic lower bound after shrinking]
\label{A:spread}
Let $T$ be a noetherian integral scheme with fraction field $K$, and suppose $n=\cd_\ell(K)<\infty$. Every nonempty open $W\subseteq T$ contains a nonempty affine open $U$ such that
\[
\cd_\ell(U')\geq n
\qquad\text{for every nonempty open }U'\subseteq U.
\]
No excellence assumption is needed.
\end{proposition}
\begin{proof}
For $n=0$ the assertion is immediate. For $n>0$, the definition of cohomological dimension and filtered-colimit compatibility give a finite discrete $\mathbf F_\ell[\Gal(K^{\mathrm{sep}}/K)]$-module $M$ and a nonzero class $c\in H^n(K,M)$. Here we use the standard reduction to $\ell$-torsion, so that the module can be chosen finite. The finite Galois action defining $M$ spreads to a locally constant constructible sheaf $\mathcal M$ on a nonempty affine open of $W$. Continuity for the inverse system of nonempty affine opens, whose limit is $\Spec K$, then spreads $c$ to $c_U\in H^n(U,\mathcal M)$. If its restriction to a nonempty open $U'$ vanished, its further restriction to $\Spec K$ would vanish, contradicting the choice of $c$. Thus the same sheaf and class witness the claimed lower bound on every such $U'$. We use the precise continuity theorem \citet[Exposé VII, Théorème 5.7]{A_SGA4}, also used in \citet[Exposé XVIII$_A$, proof of Corollary 1.2]{A_ILO2014}.
\end{proof}

\subsection{A.1: an excellent normal counterexample}

\begin{theorem}[Refuted; A.1 and the proposed global passage in L1]
\label{A:realcone}
The assertion on 31/085 that an integral affine noetherian scheme $T$ with $\ell$ invertible satisfies $\cd_\ell(T)\leq\cd_\ell(\Frac T)$ is false even when $T$ is excellent and normal. More precisely, for
\[
T=\Spec\bigl(\R[x,y,z]/(x^2+y^2+z^2)\bigr),\qquad \ell=2,
\]
we have $\cd_2(T)=+\infty$ and $\cd_2(\Frac T)=2$.
\end{theorem}
\begin{proof}
The quadratic form $x^2+y^2+z^2$ is irreducible over $\R$, so the displayed ring is a domain. It is a finitely generated algebra over a field and hence excellent. Its dimension is two. A hypersurface ring is Cohen--Macaulay, hence satisfies Serre's condition $S_2$. Its singular locus is the origin, since the three partial derivatives are $2x,2y,2z$. The origin has codimension two. Thus the ring is regular in codimension one and is normal by Serre's criterion. The prime $2$ is a unit.

Let $K$ be its fraction field. In $K$ the element $x$ is nonzero, and
\[
-1=(y/x)^2+(z/x)^2.
\]
Consequently $K$ is not orderable. The algebraic extension $K(i)$ is a function field of transcendence degree two over the algebraically closed field $\C$, so $\cd_2(K(i))=2$. For a nonorderable field, finite algebraic extension preserves its $2$-cohomological dimension; hence $\cd_2(K)=2$. These field assertions are \citet[Exposé X, Théorème 2.1 and its finite-extension input]{A_SGA4}.

The origin defines a closed immersion $i:\Spec\R\hookrightarrow T$. For $\Lambda=\Z/2$, exactness of $i_*$ gives
\[
H^q(T,i_*\Lambda)=H^q(\R,\Lambda)=H^q(C_2,\F_2)=\F_2
\quad(q\geq0).
\]
The last equality follows from the periodic resolution of the cyclic group of order two: with trivial $\F_2$ coefficients, both alternating differentials vanish. Thus no finite bound is possible. The same real quadric obstruction is explicitly noted in \citet[Exposé X, Exemple 4.5]{A_SGA4}; our verification includes its normality and the exact two dimensions needed here.
\end{proof}

\begin{proposition}[Known; a valid general replacement]
\label{A:general-upper}
For a finite-dimensional noetherian scheme $T$ and a prime $\ell$ invertible on $T$,
\[
\cd_\ell(T)\leq
\sup_{t\in T}\bigl(\cd_\ell(\kappa(t))+2\dim\OO_{T,t}\bigr).
\]
This estimate does not require quasi-excellence or affineness. Unlike the assertion in A.1, it retains the cohomological dimensions of nongeneric residue fields.
\end{proposition}
\begin{proof}
This is exactly \citet[Exposé XVIII$_A$, Lemma 2.2]{A_ILO2014}, with the same convention of all torsion sheaves. The estimate is vacuous for the real quadric because the residue field at the vertex has infinite $2$-cohomological dimension. It therefore makes no invalid passage from the generic point to its specializations.
\end{proof}

\begin{remark}[Scope of the refutation]
The counterexample is quasi-excellent and hence also refutes A.1 without that hypothesis. It does not refute the eventual equality asked in A.2: one can remove the offending closed point. Nor does it refute a statement in which all relevant residue fields are assumed nonorderable. Those distinctions matter when evaluating the proposed fibration argument in L1.
\end{remark}


\begin{corollary}[Known after Gabber; a generic-field replacement without excellence]
\label{A:generic-plus-dimension}
Let $X$ be a finite-dimensional noetherian integral scheme, let $\ell$ be invertible on $X$, and put $K=\Frac(X)$. If $\ell=2$, assume that no residue field of $X$ is orderable. Then
\[
 \cd_\ell(X)\leq\cd_\ell(K)+\dim X\leq2\cd_\ell(K).
\]
The first estimate is sharp even for excellent regular affine schemes.
\end{corollary}
\begin{proof}
Theorem~\ref{A:strict-field} supplies, without excellence, the local fraction-field hypothesis in \citet[Exposé X, Corollaire 4.4]{A_SGA4}. That corollary gives the first inequality. Its proof applies the local specialization lower bound and the dimension-function criterion; no affine hypothesis is used. The second inequality is \citet[Exposé X, Corollaire 2.5]{A_SGA4}. The excellent Dedekind scheme of Theorem~\ref{A:nagata} has dimensions $3$, $2$, and $1$, respectively, and therefore realizes equality in the first bound.
\end{proof}

\subsection{A.2: a proved eventual-equality case beyond separably closed fields}

\begin{lemma}[Proved; quadratic descent with a finite dimension bound]
\label{A:quadratic-descent}
Let $f:T'\to T$ be a torsor under the constant group $\Z/2$, and suppose $\cd_2(T)<\infty$. Then $\cd_2(T)=\cd_2(T')$.
\end{lemma}
\begin{proof}
Finite étale pushforward is exact, so $\cd_2(T')\leq\cd_2(T)$. Put $a=\cd_2(T')$. For every sheaf $F$ of $\F_2$-vector spaces on $T$ there is an exact sequence
\[
 0\longrightarrow F\longrightarrow f_*f^*F
   \longrightarrow F\longrightarrow0.
\]
The first map is the diagonal and the second is the sum of the two coordinates, after an étale trivialization of the torsor. The sequence therefore descends. For $q>a$, its cohomology sequence gives
$H^q(T,F)\simeq H^{q+1}(T,F)$. Since the latter groups vanish in sufficiently high degrees by hypothesis, they vanish for every $q>a$. Reduction to $\F_2$ coefficients proves $\cd_2(T)\leq a$.
\end{proof}

\begin{theorem}[Proved under stated hypotheses; A.2 over fields]
\label{A:eventual-field}
Let $X$ be an integral affine scheme essentially of finite type over a field $k$, and let $\ell\neq\operatorname{char}k$. Suppose $n=\cd_\ell(\Frac X)<\infty$. Every nonempty open $W$ of $X$ contains a nonempty affine open $U$ such that
\[
 \cd_\ell(U')=n
 \quad\hbox{for every nonempty affine open }U'\subset U.
\]
If $\cd_\ell(k)<\infty$, the upper bound $\cd_\ell(V)\leq n$ holds for every affine open $V$ of $X$, before shrinking.
\end{theorem}
\begin{proof}
First suppose $X$ is of finite type, and put $d=\trdeg_k\Frac X$. If $\cd_\ell(k)<\infty$, Artin's affine theorem \citep[Exposé XIV, Corollaire 3.2]{A_SGA4}, followed by Hochschild--Serre descent, gives
\[
 \cd_\ell(V)\leq d+\cd_\ell(k)=n
\]
for every affine open $V$. The equality on the right is \citet[Exposé X, Théorème 2.1]{A_SGA4}.

That same field theorem says that the only remaining possibility is $\ell=2$, $k$ orderable, and $K=\Frac X$ nonorderable. In this case $\cd_2(k(i))<\infty$ and
\[
 n=d+\cd_2(k(i)).
\]
Write $-1$ as a finite sum of squares in $K$ and invert a common nonzero denominator. On the resulting nonempty open $X_1$, the identity remains valid in every residue field, so none is orderable. For every affine open $V\subset X_1$, Corollary~\ref{A:generic-plus-dimension} gives finite $\cd_2(V)$. The quadratic étale cover $V_{k(i)}\to V$ and Lemma~\ref{A:quadratic-descent} now give
\[
 \cd_2(V)=\cd_2(V_{k(i)})\leq d+\cd_2(k(i))=n.
\]
This also covers the case that the quadratic cover is split.

For the essentially finite-type version, write the coordinate ring as a localization of a finitely generated integral $k$-algebra. Any affine open and any constructible sheaf on it descend to an affine open at a finite localization stage. In the orderable case first invert the finitely many denominators in the displayed sum of squares. The upper bounds just proved hold at every such stage; continuity \citep[Exposé VII, Théorème 5.7]{A_SGA4} carries them to the localization. Finally apply Proposition~\ref{A:spread} inside $W$, intersecting with $X_1$ when necessary. Its one surviving generic class gives the opposite inequality simultaneously on every smaller nonempty affine open.
\end{proof}

\begin{proposition}[Known with derivation; the arithmetic test in the earlier study]
\label{A:arithmetic-test}
Let $K$ be a totally imaginary number field and $X=\Spec\mathcal O_K$. For every prime $\ell$,
\[
 \cd_\ell(X)=3,\qquad
 \cd_\ell(U)=2
 \quad (\,\varnothing\neq U\subsetneq X\text{ open}\,).
\]
Consequently the study is right that $X$ is not a counterexample satisfying the invertibility hypothesis of A.6. Its conclusion that A.6 remains undecided is superseded by Corollary~\ref{A:nagata-global-consequences}.
\end{proposition}
\begin{proof}
We use Artin--Verdier duality in the form \citet[Chapter II, Theorem 3.1, Corollary 3.2, and Proposition 2.3(a)]{A_Milne2006}, for arbitrary finite constructible étale sheaves. Since $K$ has no real places, no modified real-place term occurs. Duality gives vanishing above degree three. If $j:X[1/\ell]\to X$ and $F=j_!\mu_\ell$, then
\[
 H^3(X,F)^\vee
 =\Hom_X(F,\Gm)
 =\Hom_{X[1/\ell]}(\mu_\ell,\Gm)\neq0,
\]
so the bound three is attained.

For a proper nonempty open $U$, the localization sequence is
\[
 \bigoplus_{v\notin U}H^2(K_v,F)\longrightarrow H_c^3(U,F)
 \longrightarrow H^3(U,F)\longrightarrow0.
\]
Global and local duality identify the dual of the first arrow with restriction
\[
 \Hom_U(F,\Gm)\longrightarrow
 \prod_{v\notin U}\Hom_{K_v}(F,\Gm).
\]
It is injective: a morphism to $\Gm$ which vanishes at the generic point vanishes on the normal integral scheme $U$, and extension from $K$ to any $K_v$ is faithful on morphisms of finite generic sheaves. There is at least one omitted finite place. Hence $H^3(U,F)=0$, and continuity and coefficient dévissage give $\cd_\ell(U)\leq2$ for all torsion sheaves.

For completeness, the value two is attained without assuming that a generic class automatically extends to all of $U$. The finite group $H^1(U,\F_\ell)$ has some dimension $b$ by the same arithmetic finiteness theorem. Choose $b+1$ distinct closed points, put $Z$ equal to their union, and let $j:U-Z\to U$. From
\[
 0\longrightarrow j_!\F_\ell\longrightarrow\F_\ell
 \longrightarrow(\F_\ell)_Z\longrightarrow0
\]
we obtain a nonzero cokernel of
$H^1(U,\F_\ell)\to H^1(Z,\F_\ell)\simeq\F_\ell^{b+1}$
inside $H^2(U,j_!\F_\ell)$. Thus $\cd_\ell(U)=2$.
Every rational prime fails to be a unit on $X$ itself. On $X[1/\ell]$, the number $n$ in A.6 is indeed two: the generic term is two and each closed-point term is $1+1$.
\end{proof}


\subsection{An approximation bound and a discarded Dedekind construction}

\begin{lemma}[Proved: approximation by finitely generated subalgebras]
\label{A:A2-approximation}
Let $k$ be a field, let $\ell\ne\operatorname{char}k$, and suppose
$c=\operatorname{cd}_{\ell}k<\infty$.  Let $R$ be a $k$-algebra domain
whose fraction field $K$ has finite transcendence degree $t$ over $k$.
We do not assume that $R$ is noetherian or finitely generated.  Then
\[
\operatorname{cd}_{\ell}(\operatorname{Spec}R)\le c+t.
\]
In particular, if $k$ is algebraically closed, $K$ is algebraic over
$k(t)$, and $R$ is a Dedekind subring of $K$ with fraction field $K$,
then $R$ contains $k$, its $\ell$-cohomological dimension is at most
one, and $\operatorname{Pic}(R)/\ell=0$.
\end{lemma}

\begin{proof}
Write $R=\varinjlim_i R_i$, where $R_i$ runs through its finitely
generated $k$-subalgebras.  Every $R_i$ is a domain and
$\dim R_i=\operatorname{trdeg}_k\operatorname{Frac}(R_i)\le t$.
Artin's affine theorem over a separable closure of $k$, followed by
the Hochschild--Serre spectral sequence, gives
$\operatorname{cd}_{\ell}(\operatorname{Spec}R_i)\le c+t$; the
precise affine theorem is SGA~4, Expos\'e~XIV, Corollary~3.2
\cite{C_SGA4}.

Here are the two limit steps, including nonconstructible coefficients.
Every constructible $\mathbf Z/\ell^n$-sheaf on
$\operatorname{Spec}R$ descends to some $\operatorname{Spec}R_i$,
and its cohomology is the filtered colimit of the cohomology of its
pullbacks to later stages.  Consequently it has no cohomology in
degrees above $c+t$.  Every $\ell$-primary torsion sheaf is a filtered
colimit of constructible $\ell$-primary torsion sheaves, and cohomology
on the affine scheme $\operatorname{Spec}R$ commutes with that
filtered colimit.  We use, respectively, the Stacks Project,
Lemma~59.73.10(3), Theorem~59.51.3, Lemma~59.73.2(2)--(3), and
Lemma~59.51.4, all opened in this run \cite{C_StacksLimits}.
This proves the assertion for arbitrary torsion coefficients.

For the final assertion, consider a maximal ideal $\mathfrak m$ of
the Dedekind domain $R$ and its normalized discrete valuation
$v_{\mathfrak m}\colon K^\times\to\mathbf Z$.  The group $k^\times$
is divisible because $k$ is algebraically closed.  Its image in
$\mathbf Z$ is therefore zero.  Thus every nonzero element of $k$
is a unit in every $R_{\mathfrak m}$.  The intersection identity
$R=\bigcap_{\mathfrak m}R_{\mathfrak m}$ inside $K$ gives $k\subset R$.
If $R$ is itself a field, this inclusion is immediate.  We may now
apply the first assertion with $c=0$ and $t=1$.  Finally the Kummer
sequence injects $\operatorname{Pic}(R)/\ell$ into
$H^2(R,\mu_\ell)=0$.
\end{proof}

\begin{proposition}[Proved; the algebraic-number-field case of generic dimension two]
\label{A:A2-arithmetic-two}
Let $R$ be a characteristic-zero domain whose fraction field $K$ is
algebraic over $\Q$, and let $p$ be invertible in $R$. If $p=2$,
suppose that $K$ is not orderable. Then
\[
 \cd_p(\Spec R)\leq2.
\]
No noetherian hypothesis is needed for this upper bound. If $R$ is
noetherian and $\cd_p(K)=2$, then A.2 holds for $\Spec R$.
\end{proposition}
\begin{proof}
Write $R$ as the filtered union of its finitely generated
$\Z[1/p]$-subalgebras $R_i$. Each $F_i=\Frac(R_i)$ is a number
field. If $p=2$, choose one identity
$-1=\sum_j (a_j/b_j)^2$ in $K$, with $a_j,b_j\in R$ and
$b_j\ne0$, and restrict to the cofinal system containing all these
elements. Every resulting $F_i$ is then totally imaginary.

The finite-type one-dimensional domain $R_i$ has finite
normalization $\widetilde R_i$. Its spectrum is an open subscheme
of $\Spec\mathcal O_{F_i}$: a normal finite-type model of a number
field over $\Spec\Z$ is quasi-finite, and the birational form of
Zariski's main theorem identifies it with such an open. This open
omits the primes over $p$, because $p$ is a unit. Thus it is a
proper open of the arithmetic curve. Artin--Verdier duality in
\citet[Chapter II, Theorem 3.1, Corollary 3.2, and Proposition 2.3(a)]{A_Milne2006},
in the form verified in Proposition~\ref{A:arithmetic-test}, gives
$\cd_p(\widetilde R_i)\leq2$. For odd $p$ the real-place terms have
no positive $p$-cohomology; for $p=2$ there are no real places in
the chosen system.

Let $\nu:\Spec\widetilde R_i\to\Spec R_i$ be the finite
normalization, and let $\mathcal F$ be a constructible sheaf of
$\mathbf F_p$-vector spaces on $\Spec R_i$. The adjunction map
\[
 0\longrightarrow\mathcal F\longrightarrow
 \nu_*\nu^*\mathcal F\longrightarrow\mathcal Q\longrightarrow0
\]
is injective: on each geometric stalk it is the diagonal map into
the nonempty finite collection of stalks over that point. The
cokernel is supported on the finite closed set where $\nu$ is not
an isomorphism. Every residue field there is finite, so
$H^q(\Spec R_i,\mathcal Q)=0$ for $q>1$. Finite pushforward is
exact and preserves cohomology. The long exact sequence therefore
gives $H^q(\Spec R_i,\mathcal F)=0$ for $q>2$.
Coefficient d\'evissage gives the same bound for all $p$-primary
torsion sheaves.

Constructible descent, cohomological continuity, and approximation
of arbitrary torsion sheaves by constructible ones, with the exact
locators in Lemma~\ref{A:A2-approximation}, pass this uniform bound
to $R$. Every affine open of $\Spec R$ has the same type of
coordinate domain and the same fraction field. If its generic
dimension is two, Proposition~\ref{A:spread} supplies the matching
lower bound after shrinking, proving A.2.
\end{proof}

\begin{corollary}[Proved; why the completely split positive-$\mu$ towers do not refute A.2]
\label{A:A2-split-tower-exclusion}
Let $K_\infty/F$ be a $\Z_p$-extension of a number field, with $F$
totally imaginary if $p=2$. Suppose that some finite prime of $F$
not dividing $p$ splits completely in $K_\infty/F$. Then
$\cd_p(K_\infty)=2$. Every noetherian subring $R\subset K_\infty$
with fraction field $K_\infty$ and with $p$ invertible satisfies A.2.
The integral closure of $\mathcal O_F[1/p]$ in $K_\infty$ is not
noetherian.
\end{corollary}
\begin{proof}
Choose a compatible prime over the completely split prime.
The union of the corresponding discrete valuation rings is a
discrete valuation ring: all ramification indices and residue
degrees are one. Its residue field is the original finite field.
Lemma~\ref{A:valuation-lower} therefore gives
$\cd_p(K_\infty)\geq2$. The opposite inequality follows from
$\cd_p(F)=2$ and monotonicity under algebraic extension, with the
real-place convention as stated. Proposition~\ref{A:A2-arithmetic-two}
now applies to every such $R$ and to all its affine opens.

The integral closure has infinitely many distinct primes over the
chosen prime, since their number in the finite layers tends to
infinity. It is a one-dimensional domain integral over a Dedekind
domain. A nonzero element of the chosen prime of $F$ lies in all
these infinitely many maximal ideals. In a noetherian
one-dimensional domain the primes containing a given nonzero
element form a finite set, since its quotient is a zero-dimensional
noetherian ring. This is a contradiction.

The split-prime hypothesis is present in the constructions checked
in \citet[Proposition 3.1 and Corollary 3.2, pp.~11--12]{A_HajirMaire}
and \citet[Lemma 1 and \S3, Theorem 2]{A_HubbardWashington}.
The conclusion does not depend on their positive-$\mu$ calculation:
the split valuation alone forces generic dimension two, and
localizing the ring does not change that fraction field.
\end{proof}

\begin{lemma}[Proved; why a pseudo-algebraically closed fraction field cannot supply the proposed counterexample]
\label{A:A2-PAC-exclusion}
A pseudo-algebraically closed field admits no nontrivial discrete
valuation. Consequently it cannot be the fraction field of a
positive-dimensional noetherian domain.
\end{lemma}
\begin{proof}
Suppose $v:K^\times\to\Z$ is a surjective discrete valuation, and
choose $t\in K$ with $v(t)=1$. Choose a prime $q$ different from
$\operatorname{char}K$, and an integer $m\geq2$ also prime to that
characteristic. The affine curve
\[
 C:\quad y^q=t(1+t^q x^{mq})
\]
is geometrically integral. Indeed $1+t^q x^{mq}$ has simple nonzero
roots over $\overline K$. At any one of them the displayed monic
polynomial in $y$ is Eisenstein, hence irreducible over
$\overline K(x)$ and, by Gauss's lemma, over $\overline K[x,y]$.

It has no $K$-point. If $v(x)\geq0$, the valuation of its right-hand
side is one. If $v(x)<0$, it is
$1+q+mqv(x)$, again congruent to one modulo $q$. Neither is the
valuation of a nonzero $q$th power. The case $x=0$ belongs to the
first alternative, and the right-hand side is never zero for a
valued $K$-point. This contradicts the defining rational-point
property of a pseudo-algebraically closed field.

If a noetherian domain $R$ of positive dimension had fraction field
$K$, a height-one localization and then a localization of its
normalization would give a discrete valuation ring with fraction
field $K$. The normalization step uses the one-dimensional theorem
already opened at \citet[Tag 09IG, Lemma 10.120.18]{A_Stacks}; it
does not require finite normalization or excellence. This proves
the final assertion. The exclusion is classical; we include the
elementary curve argument to verify the failed construction route
without relying on an unopened textbook locator.
\end{proof}

Lemma~\ref{A:A2-approximation} rules out a persistent Picard-group obstruction obtained by
prescribing a Dedekind class group while keeping the fraction field
algebraic over a one-variable function field over an algebraically
closed field.  It does not rule out algebraic extensions of number
fields: there the cohomological dimension of the fraction field can
drop below the bound furnished by the original ground field.

\begin{corollary}[Proved under stated hypotheses; an additional A.2 case]
\label{A:A2-approximation-equality}
Let $R$ be a noetherian $k$-algebra domain, with $p\ne\operatorname{char}k$,
and suppose
\[
 \cd_p(\Frac R)=\cd_p(k)+\operatorname{trdeg}_k\Frac R<\infty.
\]
Then A.2 holds for $\Spec R$, without a finite-type or excellence
assumption on $R$.
\end{corollary}
\begin{proof}
Every affine open is the spectrum of a $k$-algebra domain with the
same fraction field. Lemma~\ref{A:A2-approximation} gives the upper
bound by the displayed common value on every such open.
Proposition~\ref{A:spread} gives the matching lower bound after
shrinking inside any prescribed nonempty open, and on every smaller
nonempty affine open.
\end{proof}

\begin{lemma}[Proved: finite character in an algebraic tower]
\label{A:A2-tower-Dedekind}
Let $R_0$ be a Dedekind domain and let $K_\infty/K_0$, with
$K_0=\operatorname{Frac}(R_0)$, be an algebraic extension expressed as
a directed union of finite extensions $K_i$.  Suppose that their
integral closures $R_i$ of $R_0$ are finite over $R_0$.  Suppose also
that each nonzero prime of every $R_i$ has only finitely many primes
above it in $R_\infty=\bigcup_iR_i$, and that along each such prime
the ramification indices in the tower are bounded.  Then
$R_\infty$ is a Dedekind domain.
\end{lemma}

\begin{proof}
Integrality shows that every nonzero prime of $R_\infty$ is maximal.
At such a prime $P$, the localization $(R_\infty)_P$ is the directed
union of the compatible valuation rings $(R_i)_{P\cap R_i}$.  It is
a valuation ring.  Bounded ramification indices make its value group
a discrete cyclic subgroup of $\mathbf Q$, so this localization is
a discrete valuation ring.  A nonzero element $a\in R_\infty$
belongs to some $R_i$ and has only finitely many prime divisors in
that $R_i$; the finite-splitting hypothesis then gives only finitely
many maximal ideals of $R_\infty$ containing $a$.

We check noetherianity directly.  Given a nonzero ideal $I$, choose
$0\ne a\in I$, and list the finitely many maximal ideals
$P_1,\dots,P_s$ containing $a$.  For each $j$ choose $b_j\in I$
whose valuation at $P_j$ is the minimum of the nonnegative integer
valuations attained on $I\setminus\{0\}$.  The ideal
$J=(a,b_1,\dots,b_s)$ has $J_{P_j}=I_{P_j}$ for every $j$; at every
other maximal ideal $a$ is a unit, so again $J_P=I_P$.  Local equality
of these ideals at all maximal ideals implies $I=J$.  Thus every
ideal is finitely generated.  The domain is noetherian, normal and
one-dimensional, as required.
\end{proof}

\subsection{A.2: an exact cyclotomic arithmetic criterion}

We now isolate an arithmetic instance of the eventual upper bound. This is
a proved equivalence, not a proof of either the general conjecture in A.2
or Iwasawa's $\mu$-conjecture. We use $\mu(L)=0$ to mean vanishing of the
class-group Iwasawa invariant of the cyclotomic $\Z_p$-field $L$; the
vanishing condition is independent of the finite initial layer.

\begin{lemma}[Proved from known arithmetic inputs; a cyclotomic Dedekind domain]
\label{A:cyclotomic-domain}
Fix a prime $p$ and a number field $F$ containing $\mu_p$; if $p=2$,
assume $i\in F$. Let $K=F_{\mathrm{cyc}}$ be its cyclotomic
$\Z_p$-extension, and put $R=\mathcal O_K[1/p]$. Then $R$ is an excellent
Dedekind domain, $\cd_p(K)=1$, and every residue field at a closed point
of $\Spec R$ has $p$-cohomological dimension zero. The same assertions
hold after replacing $K$ by any finite extension $L/K$ and $R$ by its
normalization $R_L=\mathcal O_L[1/p]$.
\end{lemma}
\begin{proof}
Write $K=\bigcup_n F_n$ with $[F_n:F]=p^n$. At a finite prime
$v\nmid p$ of $F$, the tower is unramified. Its Frobenius has infinite
order in $\operatorname{Gal}(K/F)\simeq\Z_p$: under the cyclotomic
character, a zero image would say that the positive integer $Nv>1$ is
a $p$-adic root of unity. Hence the decomposition group is open, so
only finitely many primes lie above $v$. The ramification indices are
one. Lemma~\ref{A:A2-tower-Dedekind} therefore proves that $R$ is
Dedekind. The cyclotomic ramification statements used here are also
\citet[Proposition 11.1.1]{A_NSW}; at primes above $p$ the inertia
subgroups are open and the tower is eventually totally ramified.
A Dedekind domain whose fraction field has characteristic zero is
excellent by \citet[Proposition 15.53.3(4), Tag 07QS]{A_Stacks}.

Every finite extension $L/K$ has the form $F'K$ for a number field
$F'$ containing $F$. Since $K=F\Q_{\mathrm{cyc}}$, this is exactly
the cyclotomic $\Z_p$-extension of $F'$, after changing its finite
initial layer. Repeating the preceding argument proves the assertions
about $R_L$; excellence also makes $R_L$ finite over $R$.

We check the fraction-field dimension, including the passage to arbitrary
coefficients. For every such $L$, a class of $\operatorname{Br}(L)[p]$
descends to a finite layer. Brauer--Hasse--Noether injects the Brauer
group of that number field into the direct sum of its local Brauer
groups. Only finitely many local invariants occur, and restriction to
a finite local extension multiplies an invariant by the local degree.
At every finite place these local degrees in the cyclotomic tower are
unbounded powers of $p$: at $p$ by ramification and elsewhere by the
nonzero Frobenius just checked. Thus one later layer kills every local
invariant simultaneously. There are no real places. Consequently
\[
 \operatorname{Br}(L)[p]=0
 \quad\hbox{for every finite }L/K.
\]
The precise local-global and degree inputs, including the unbounded
local degrees, are \citet[Chapter VII, Theorem 7.1, Proposition 7.2,
and the proof of Lemma 7.3]{A_MilneCFT}.

Since $K$ is an algebraic extension of the totally imaginary number
field $F$, initially $\cd_p(K)\leq2$ by \citet[Proposition 8.3.18, taking all primes in $S$]{A_NSW} and closed-subgroup monotonicity. If $M$ is a finite discrete
$\F_p[\operatorname{Gal}(K^{\mathrm{sep}}/K)]$-module, choose a finite
Galois extension $L/K$ trivializing it. A sum of regular permutation
modules surjects onto $M$. Its second cohomology vanishes by Shapiro,
Kummer theory, and $\mu_p\subset L$, because
$H^2(L,\F_p)\simeq\operatorname{Br}(L)[p]=0$.
The kernel has no third cohomology by the preliminary bound two.
The long exact sequence gives $H^2(K,M)=0$. Filtered colimits and
coefficient d\'evissage give $\cd_p(K)\leq1$. To see equality, choose
a prime above $v\nmid p$ and an element of $F$ of valuation one at $v$.
Its valuation remains one in $K$, so it is not a $p$th power. Its Kummer
class is nonzero in $H^1(K,\mu_p)$.

Finally, a residue field of $R$ is the union of finite fields whose
successive degrees are eventually $p$. Its absolute Galois group has
no nontrivial pro-$p$ subgroup, so its $p$-cohomological dimension is
zero. This also follows from the valuation inequality
$1+\cd_p\kappa(v)\leq\cd_p(K)=1$ already proved in Part A.
\end{proof}

\begin{lemma}[Proved from Iwasawa's structure theorem; the persistent obstruction]
\label{A:cyclotomic-Pic}
Use the notation of Lemma~\ref{A:cyclotomic-domain}. For every nonempty
open $U\subset\Spec R$,
\[
 H^2(U,\mu_p)\simeq\operatorname{Pic}(U)/p.
\]
If $\mu(K)=0$, these groups are zero. If $\mu(K)>0$, they have infinite
dimension over $\F_p$, even after deleting any finite set of closed
points. The assertions hold for every finite extension $L/K$ as well.
\end{lemma}
\begin{proof}
Let $C_K=\varinjlim_n\operatorname{Cl}(\mathcal O_{F_n})$, where the
maps are extension of ideals, and let $A_K$ be its $p$-primary part.
The Iwasawa class-group structure theorem gives the following two
facts about $M=A_K^\vee$, its Pontryagin dual: $M$ is a finitely
generated torsion $\Lambda=\Z_p[[T]]$-module with the class-group
$\mu$-invariant, and it has no nonzero finite $\Lambda$-submodule.
We use \citet[Proposition 3.2.15 and Corollary 3.2.16]{A_SharifiIwasawa},
whose proofs were opened, together with the structure theorem in
Theorem 2.2.13 of that source.

If $\mu(K)=0$, the structure theorem makes $M$ finitely generated
over $\Z_p$. Its $p$-power torsion is then finite, hence zero by the
absence of finite $\Lambda$-submodules. Thus $M$ is free over $\Z_p$
and $A_K$ is $p$-divisible. If $\mu(K)>0$, the elementary module of
$M$ has a summand $\Lambda/(p^a)$ with $a\geq1$. Its subgroup
annihilated by $p$ is infinite. A pseudo-isomorphism changes this
subgroup only by finite kernels and cokernels, so $M[p]$ is infinite.
Duality identifies $M[p]=(A_K/pA_K)^\vee$, proving that $A_K/pA_K$
has infinite $\F_p$-dimension. The prime-to-$p$ part of $C_K$ is
$p$-divisible, hence the same alternatives hold for $C_K/pC_K$.

Invertible modules descend through a filtered union of rings, as do their isomorphisms. Hence the Picard group of $R=\varinjlim_n\mathcal O_{F_n}[1/p]$ is the direct limit of the finite-layer Picard groups. Localization of ideals therefore gives
\[
 \operatorname{Pic}(R)=C_K/D_p,
\]
where $D_p$ is the image of the direct limit of the groups generated
by primes above $p$. After a finite layer these primes no longer
split and are totally ramified at each step. That direct limit is a
finite direct sum of copies of $\Z[1/p]$, so $D_p=pD_p$.
In particular $D_p\subset pC_K$, and
\[
 \operatorname{Pic}(R)/p\simeq C_K/pC_K.
\]
Every nonempty open in the noetherian one-dimensional scheme
$\Spec R$ has finite complement. Its Picard group is the quotient
of $\operatorname{Pic}(R)$ by the classes of those finitely many
primes. Such a quotient preserves vanishing modulo $p$, and it
cannot kill an infinite-dimensional quotient modulo $p$.

For clarity, the Brauer comparison used in Kummer theory also applies
to these possibly non-finite-type Dedekind schemes. For the generic
inclusion $j:\Spec K\to U$, the divisor sequence of étale sheaves is
\[
 0\longrightarrow\Gm\longrightarrow j_*\Gm
   \longrightarrow\bigoplus_{x\in U^{(1)}}i_{x*}\Z\longrightarrow0.
\]
At strictly henselian localizations its exactness is the valuation
sequence of a discrete valuation ring. Hilbert 90 gives
$R^1j_*\Gm=0$, while $H^1(\kappa(x),\Z)=0$. The first cohomology of the displayed direct sum is zero as well: express it as the filtered union of finite sums and use cohomological continuity on the quasi-compact, quasi-separated scheme $U$. The long exact sequence
and the Leray sequence therefore inject
$H^2(U,\Gm)$ into $H^2(K,\Gm)=\operatorname{Br}(K)$.
The latter has no $p$-torsion by Lemma~\ref{A:cyclotomic-domain}.
Kummer theory now gives the asserted isomorphism. The proof for
$L/K$ is identical.
\end{proof}

\begin{theorem}[Proved; A.2 on cyclotomic Dedekind domains is equivalent to a precise arithmetic condition]
\label{A:Iwasawa-equivalence}
For $R$ and $K$ in Lemma~\ref{A:cyclotomic-domain}, the following
conditions are equivalent:
\begin{enumerate}[label=(\roman*)]
\item Some nonempty affine open $U\subset\Spec R$ has $\cd_p(U)\leq1$.
\item Every nonempty affine open $U\subset\Spec R$ has $\cd_p(U)=1$.
\item For every finite extension $L/K$, the cyclotomic class-group
      invariant satisfies $\mu(L)=0$.
\end{enumerate}
Consequently A.2 for this excellent one-dimensional domain holds
if and only if (iii) holds. If (iii) fails, every nonempty affine open
of $\Spec R$ has $p$-cohomological dimension exactly two.
\end{theorem}
\begin{proof}
Let $U$ be any nonempty affine open. Corollary~\ref{A:generic-plus-dimension}
gives the preliminary bound $\cd_p(U)\leq2$. Its closed residue fields
have $p$-cohomological dimension zero.

Suppose (i), and let $f:Y\to U$ be the finite normalization in
a finite extension $L/K$. It is an open arithmetic Dedekind scheme
of the kind in Lemma~\ref{A:cyclotomic-Pic}. Finite pushforward of
torsion étale sheaves is exact and computes cohomology:
\[
 H^q(Y,E)=H^q(U,f_*E).
\]
Thus $\cd_p(Y)\leq\cd_p(U)\leq1$.
If $\mu(L)>0$, Lemma~\ref{A:cyclotomic-Pic} contradicts this with
$E=\mu_p$ and $q=2$. This proves (iii).

Conversely assume (iii). Let $F_0$ be a constructible sheaf of
$\F_p$-vector spaces on $U$. Choose a finite Galois extension $L/K$
trivializing its generic fibre, let $f:Y\to U$ be the normalization,
and choose a dense open $j:V\hookrightarrow U$ on which $F_0$ is
locally constant and $f$ is étale. Write
$j':V_Y=f^{-1}(V)\hookrightarrow Y$. A surjection from a finite sum
of regular permutation modules on the generic fibre spreads, after
shrinking $V$ if necessary, to
\[
 (f|_{V_Y})_*\F_p^r\twoheadrightarrow F_0|_V.
\]
Extension by zero and finite pushforward give a morphism
\[
 P=f_*j'_!\F_p^r=j_!(f|_{V_Y})_*\F_p^r\longrightarrow F_0
\]
whose cokernel $Q$ is supported on finitely many closed points.
The equality of the two sheaves is checked on geometric stalks,
using finiteness of $f$; it is also proper base change for the open
extension-by-zero square. Let $I$ be its image and $N$ its kernel.
The preliminary dimension bound gives $H^3(U,N)=0$, and
$H^q(U,Q)=0$ for $q>0$. The two short exact sequences for $I$
therefore make $H^2(U,F_0)$ a quotient of $H^2(U,P)$.

On $Y$, the complement of $V_Y$ is finite and has residue fields
of $p$-cohomological dimension zero. The localization sequence
therefore gives
\[
 H^2(U,P)=H^2(Y,j'_!\F_p)^r
          \simeq H^2(Y,\F_p)^r=0.
\]
The last equality is Lemma~\ref{A:cyclotomic-Pic}, (iii), and
$\mu_p\simeq\F_p$ on $Y$. Hence $H^2(U,F_0)=0$ for every
constructible $F_0$. The preliminary bound, filtered-colimit
continuity and coefficient d\'evissage give $\cd_p(U)\leq1$ for
arbitrary $p$-primary torsion sheaves.

There are infinitely many closed points on $U$. Choose two,
write $Z$ for their union and $j_0:U-Z\hookrightarrow U$.
Since $U$ is connected, the cokernel of
\[
 H^0(U,\F_p)=\F_p\longrightarrow H^0(Z,\F_p)=\F_p^2
\]
is nonzero and injects into $H^1(U,j_{0!}\F_p)$. Thus $\cd_p(U)=1$,
proving (ii); (ii) implies (i) immediately.

If (iii) fails, the finite-pushforward obstruction applies to every
nonempty affine $U$, whereas the preliminary bound remains two.
Finally $\cd_p(K)=1$, so the required eventual equality of A.2
is precisely equivalent to (i).
\end{proof}

\begin{corollary}[Proved; an exact Picard criterion for the one-dimensional gate]
\label{A:Dedekind-Pic-criterion}
Let $R$ be an excellent Dedekind domain, let $p$ be invertible in $R$,
and suppose $\cd_p(\Frac R)=1$. For any nonempty affine open
$U\subset\Spec R$, the following are equivalent:
\begin{enumerate}[label=(\roman*)]
\item $\cd_p(U)\leq1$;
\item $\operatorname{Pic}(Y)/p=0$ for the normalization $Y$ of $U$
in every finite extension of $\Frac R$.
\end{enumerate}
Consequently A.2 for this $R$ is equivalent to existence of one
$0\ne a\in R$ such that
\[
 \operatorname{Pic}(R_L[1/a])/p=0
 \quad\hbox{for every finite }L/\Frac R,
\]
where $R_L$ is the integral closure of $R$ in $L$.
\end{corollary}
\begin{proof}
Finite extensions have $p$-cohomological dimension at most one.
The divisor and Kummer argument in Lemma~\ref{A:cyclotomic-Pic}
therefore applies to every $Y$ here, giving
$H^2(Y,\mu_p)=\operatorname{Pic}(Y)/p$.
Normalization is finite by excellence, so exact finite pushforward
proves (i)$\Rightarrow$(ii).

For the converse, there are no orderable residue fields when $p=2$:
the regular-local case of Theorem~\ref{A:realexception}, applied to
the discrete valuation ring at such a point, would make its fraction
field orderable and its $2$-cohomological dimension infinite.
The valuation lower bound now gives $\cd_p\kappa(x)=0$ at every
closed point of $U$ and of $Y$. Corollary~\ref{A:generic-plus-dimension}
gives $\cd_p(U)\leq2$. In the permutation-sheaf resolution in the
proof of Theorem~\ref{A:Iwasawa-equivalence}, enlarge $L$ also to
contain $\mu_p$. All its comparisons now apply verbatim: finite
normalization, vanishing of positive boundary cohomology, and
vanishing of third cohomology of the kernel. Hypothesis (ii) gives
$H^2(Y,\F_p)=H^2(Y,\mu_p)=0$, which forces $H^2(U,F_0)=0$
for every constructible $\F_p$-sheaf $F_0$. Continuity and
d\'evissage prove (i).

Every nonempty open of an affine integral scheme contains a nonempty
principal open $D(a)$. Once (ii) holds on $D(a)$, it holds on every
smaller open: normalization commutes with restriction, and the
Picard group of an open Dedekind subscheme is a quotient by the
classes of its omitted primes. Hence all smaller affine opens have
dimension at most one. Proposition~\ref{A:spread} supplies the
opposite bound after one further shrink, simultaneously on all
smaller nonempty affine opens. This is exactly A.2. Its converse
implies (i) on a principal subopen and then (ii), proving the final
equivalence.
\end{proof}

\begin{corollary}[Proved; a single arithmetic instance contains the full odd-prime mu-conjecture]
\label{A:canonical-Iwasawa}
For an odd prime $p$, put
\[
 R_p=\Z[\mu_{p^\infty},1/p],\qquad K_p=\Q(\mu_{p^\infty}).
\]
A.2 for $\Spec R_p$ at $p$ is equivalent to Iwasawa's assertion
$\mu(F_{\mathrm{cyc}}/F)=0$ for every number field $F$.
\end{corollary}
\begin{proof}
The forward implication of Theorem~\ref{A:Iwasawa-equivalence} gives
$\mu(L)=0$ for every finite $L/K_p$. For a number field $F$, take
$L=FK_p$. This is a finite extension of $K_p$, and it is a finite
extension of $F_{\mathrm{cyc}}$ of degree dividing $p-1$.
Extension and norm on the direct limits of $p$-class groups have
composite multiplication by this degree, which is an automorphism
on a $p$-primary torsion group. The $p$-class group of
$F_{\mathrm{cyc}}$ is consequently a direct summand of that of $L$.
The latter is $p$-divisible because $\mu(L)=0$, so the former is
$p$-divisible as well. The structure argument in
Lemma~\ref{A:cyclotomic-Pic} then gives $\mu(F_{\mathrm{cyc}}/F)=0$.
Conversely, every finite $L/K_p$ is the cyclotomic $\Z_p$-extension
of a number field after choosing a finite initial layer. The asserted
vanishing for all number fields gives (iii) and hence A.2.
\end{proof}

\begin{remark}[Literature check and scope]
We searched the exact combinations ``cohomological dimension'',
``cyclotomic'', ``Dedekind'', and ``mu conjecture'', and opened
the primary sources above. The related Galois-group criterion is
\citet[Theorem 11.3.7]{A_NSW}; it concerns the maximal pro-$p$
extension with restricted ramification. Our proof spells out the
additional comparisons for arbitrary constructible sheaves on
$\Spec R$ and for finite extensions ramified at points of that
scheme. We claim no priority. The sources checked in this run do
not prove the general $\mu$-conjecture; Sharifi's Conjecture 3.2.17
states it explicitly. The 2026 abelian and elliptic-curve results
found in the search do not supply its assertion for all number fields.

This leaves that particular cyclotomic instance undecided. The
unrestricted assertion A.2 is refuted independently by the non-Galois
tower in Theorem~\ref{A:A2-persistent-arithmetic-counterexample}; no
claim about the cyclotomic $\mu$-conjecture is inferred from it.
\end{remark}

\subsection{The full Dedekind gate as a Milnor residue condition}

\begin{lemma}[Proved; the top-degree residue obstruction]
\label{A:Dedekind-Milnor-obstruction}
Let $R$ be a Dedekind domain, let $p\in R^\times$ be prime, and
put $K=\Frac R$ and $n=\cd_p(K)<\infty$. Suppose that $R$ is not a
field, so $n\geq1$. For each finite separable extension $L/K$, let
$R_L$ be the integral closure of $R$ in $L$, and put $Y_L=\Spec R_L$.
Then $R_L$ is a finite $R$-algebra and a Dedekind domain. Write
$Y_L^{(1)}$ for its closed points and $\partial_y$ for the tame
residue in Milnor $K$-theory at the discrete valuation of $y$.
There is a canonical isomorphism
\[
 H^{n+1}(Y_L,\mu_p^{\otimes n})\simeq
 \operatorname{coker}\left[
 K_n^M(L)/p\xrightarrow{(\partial_y)_y}
 \bigoplus_{y\in Y_L^{(1)}}K_{n-1}^M(\kappa(y))/p
 \right].
\]
Here $K_0^M(E)=\Z$. No excellence hypothesis is imposed.
\end{lemma}
\begin{proof}
We first justify the normalization assertion rather than assume a
Japanese hypothesis. Choose a $K$-basis $b_1,\ldots,b_d$ of $L$
consisting of elements integral over $R$, by clearing the finitely
many denominators in the equations of an arbitrary basis. If
$z\in R_L$, every $\operatorname{Tr}_{L/K}(zb_i)$ is integral over
$R$ and belongs to $K$, hence belongs to $R$ because $R$ is normal.
The nondegenerate trace pairing for the separable extension $L/K$
places $R_L$ inside the finite free $R$-lattice dual to
$\sum_i Rb_i$. A submodule of a finite module over the noetherian
ring $R$ is finite. Thus $R_L$ is finite; it is normal and
one-dimensional, so is Dedekind.

Finite separable extension gives $\cd_p(L)\leq n$. For every closed
point $y$ the discrete-valuation lower bound gives
$\cd_p\kappa(y)\leq n-1$. When $p=2$, the residue field cannot be
orderable: an ordering would extend through the regular local
ring $(R_L)_y$ to $L$, contrary to $\cd_2(L)<\infty$; this is the
regular-local case proved in Proposition~\ref{A:realexception}.

For a finite set $Z\subset Y_L^{(1)}$, localization and codimension-one
purity give the residue exact sequence. Taking the filtered colimit
over $Z$, using $\bigcap_Z(Y_L-Z)=\Spec L$ and cohomological
continuity, gives
\[
 H^n(L,\mu_p^{\otimes n})\longrightarrow
 \bigoplus_{y\in Y_L^{(1)}}
 H^{n-1}(\kappa(y),\mu_p^{\otimes(n-1)})
 \longrightarrow H^{n+1}(Y_L,\mu_p^{\otimes n})
 \longrightarrow H^{n+1}(L,\mu_p^{\otimes n})=0.
\]
The purity used here is only the discrete-valuation case. On a
strict henselization the tame inertia calculation gives
$i^!\mu_p^{\otimes n}=\mu_p^{\otimes(n-1)}[-2]$; it requires
$p$ invertible, not excellence. It is the same local calculation
underlying \citet[Expos\'e X, Th\'eor\`eme 2.2]{A_SGA4}.

The norm-residue isomorphism theorem
\citep[Theorem 6.16, p.~436]{A_Voevodsky2011} identifies the first
two groups with the displayed Milnor groups. It respects the
residue arrows: the residue of a cup product of unramified unit
classes is zero, and the residue of
$(\pi)\cup(u_1)\cup\cdots\cup(u_{n-1})$ is
$(\overline u_1)\cup\cdots\cup(\overline u_{n-1})$, for a
uniformizer $\pi$ and units $u_i$. These are precisely the defining
formulas for the tame Milnor residue; multilinearity and the
Steinberg relation reduce every symbol to these cases. The
exact sequence therefore has the stated cokernel.
\end{proof}

\begin{theorem}[Proved; an exact criterion in every generic dimension]
\label{A:Dedekind-Milnor-criterion}
With the notation of Lemma~\ref{A:Dedekind-Milnor-obstruction},
the following are equivalent:
\begin{enumerate}
\item $\cd_p(\Spec R)\leq n$.
\item For every finite separable $L/K$ containing a primitive
$p$th root of unity, the residue homomorphism
\[
 K_n^M(L)/p\longrightarrow
 \bigoplus_{y\in(\Spec R_L)^{(1)}}K_{n-1}^M(\kappa(y))/p
\]
is surjective.
\end{enumerate}
Consequently A.2 for this Dedekind domain holds if and only if there
is a nonzero $a\in R$ for which condition (2) holds simultaneously
with $R_L$ replaced by $R_L[1/a]$. For $n=1$, condition (2) is
$\operatorname{Pic}(R_L)/p=0$. This extends the finite-separable
version of Corollary~\ref{A:Dedekind-Pic-criterion} without
excellence.
\end{theorem}
\begin{proof}
Suppose (1). For the finite morphism $f:Y_L\to\Spec R$, exact
finite pushforward gives
$H^q(Y_L,\mathcal E)=H^q(\Spec R,f_*\mathcal E)$ for every torsion
sheaf $\mathcal E$. Hence $\cd_p(Y_L)\leq n$, and the preceding
lemma proves (2).

Conversely, all residue fields of $R$ are nonorderable when $p=2$,
by the regular-local ordering argument just used. The preliminary
bound of Corollary~\ref{A:generic-plus-dimension} gives
$\cd_p(R)\leq n+1$. It is enough to prove that
$H^{n+1}(R,\mathcal F)=0$ for each constructible
$\mathbf F_p$-sheaf $\mathcal F$.

Choose a finite Galois extension $L/K$ that trivializes its generic
Galois module and contains $\mu_p$. Shrink to a nonempty open
$j:V\hookrightarrow\Spec R$ on which $\mathcal F$ is locally
constant and $f:Y_L\to\Spec R$ is \'etale. Put
$V_L=f^{-1}V$, let $j':V_L\hookrightarrow Y_L$, and write
$f_V:V_L\to V$ for the restricted finite \'etale map. A finite
free module over the group algebra of $\operatorname{Gal}(L/K)$
surjects onto the finite generic module. Thus, for some finite $r$,
it gives a morphism
\[
 \mathcal P=j_!(f_V)_*\mathbf F_p^r
       =f_*j'_!\mathbf F_p^r\longrightarrow\mathcal F
\]
that is surjective on $V$. Let $\mathcal N$ be its kernel,
$\mathcal I$ its image, and $\mathcal Q$ its cokernel.
The latter is supported on finitely many closed points and has
cohomology zero above $n-1$. The two short exact sequences with
middle terms $\mathcal P$ and $\mathcal F$, together with
$H^{n+2}(R,\mathcal N)=0$, show that
$H^{n+1}(R,\mathcal F)$ is a quotient of
$H^{n+1}(R,\mathcal P)$.

The finite-support complement $Y_L-V_L$ also has cohomological
dimension at most $n-1$. Therefore extension by zero and finite
pushforward give
\[
 H^{n+1}(R,\mathcal P)
 =H^{n+1}(Y_L,j'_!\mathbf F_p^r)
 =H^{n+1}(Y_L,\mathbf F_p)^r.
\]
Since $\mu_p\subset L$ and $R_L$ is integrally closed with $p$
invertible, $\mu_p$ is a constant sheaf on $Y_L$. Condition (2)
and Lemma~\ref{A:Dedekind-Milnor-obstruction} make the last group
zero. This proves (1), first for constructible $\mathbf F_p$
coefficients and then for every $p$-primary torsion sheaf by
Lemma~\ref{A:basic}.

On a smaller open, the direct sum in (2) is obtained by omitting
the removed primes; its surjectivity follows from that on the
larger open. If one $a$ as stated exists, the upper bound holds on
every smaller affine open; Proposition~\ref{A:spread} provides the
simultaneous lower bound. Conversely, A.2 gives an affine open
with the upper bound, and a principal nonempty open inside it.
Finite normalization commutes with this localization, so (2)
holds for $R_L[1/a]$. Finally, when $n=1$, the residue map is the
valuation divisor map $L^\times/(L^\times)^p\to
\bigoplus_y\Z/p$; its cokernel is exactly
$\operatorname{Pic}(R_L)/p$.
\end{proof}

\begin{corollary}[Proved; all semilocal Dedekind domains]
\label{A:semilocal-Dedekind}
Let $R$ be a semilocal Dedekind domain, let $p$ be invertible in
$R$, and suppose $\cd_p(\Frac R)<\infty$. Then
\[
 \cd_p(\Spec R)\leq\cd_p(\Frac R),
\]
and A.2 holds. Neither henselianity nor excellence is assumed.
\end{corollary}
\begin{proof}
Every finite separable normalization $R_L$ is again semilocal
Dedekind: its maximal ideals lie over the finitely many maximal
ideals of $R$, with finite fibres. Let them be
$\mathfrak m_1,\ldots,\mathfrak m_s$. By the Chinese remainder
theorem choose $\pi_i\in R_L$ with valuation one at
$\mathfrak m_i$ and valuation zero at every other maximal ideal,
prescribing its class modulo $\mathfrak m_i^2$ and modulo the
other maximal ideals. For a symbol
$\{\overline u_1,\ldots,\overline u_{n-1}\}$ in the residue field
at $\mathfrak m_i$, lift each $\overline u_j$ to an element
$u_j\in R_L$ with value one at all other maximal ideals. These
lifts avoid every maximal ideal, and are therefore global units.
The symbol $\{\pi_i,u_1,\ldots,u_{n-1}\}\in K_n^M(L)/p$ has the
desired residue at $\mathfrak m_i$ and zero residue everywhere
else. For $n=1$ this simply says that $\pi_i$ has valuation vector
the $i$th basis vector. Symbols generate the Milnor groups,
proving surjectivity for every $L$. The preceding theorem proves
the upper bound and A.2. If $R$ is a field, both claims are immediate.
\end{proof}

\begin{theorem}[Proved; normalization removes the nonnormal curve gate]
\label{A:curve-normalization}
Let $R$ be any one-dimensional noetherian domain, let $p$ be
invertible in $R$, let $K=\Frac R$ with $n=\cd_p(K)<\infty$,
and let $B$ be the normalization of $R$ in $K$. Then A.2 holds
for $\Spec R$ if and only if it holds for $\Spec B$.
Equivalently, it holds if and only if there is a nonzero $a\in R$
such that the residue maps of
Theorem~\ref{A:Dedekind-Milnor-criterion}, for the finite separable
normalizations of $B[1/a]$, are all surjective. In particular,
A.2 holds for every semilocal one-dimensional noetherian domain.
No finiteness of $B$ over $R$ and no excellence is required.
\end{theorem}
\begin{proof}
By Krull--Akizuki in the precise form
\citet[Tag 09IG, Lemma 10.120.18]{A_Stacks}, $B$ is Dedekind and
$\nu:\Spec B\to\Spec R$ is integral and surjective, with finite
fibres and finite residue extensions. In particular $n\geq1$.
When $p=2$, first restrict to a principal nonempty open on which
a fraction-field identity $-1=\sum_i u_i^2$ is defined.
No residue field on this open is orderable. This restriction
does not change the eventual assertion A.2. At a closed point
$x$ of this open, choose a point $y$ of its normalization above
it. The valuation bound gives $\cd_p\kappa(y)\leq n-1$.
Finite-extension invariance of field cohomological dimension,
with the real exception excluded, gives
$\cd_p\kappa(x)\leq n-1$ as well.

We give the needed sheaf comparison even for an infinite
normalization. Pushforward under an integral morphism is exact
on all abelian \'etale sheaves
\citep[Tag 04C2, Lemma 59.43.5]{A_Stacks} and commutes with
arbitrary base change
\citep[Tag 0EYP, Lemma 59.55.4]{A_Stacks}. For a constructible
$\mathbf F_p$-sheaf $\mathcal F$ on $\Spec R$, the unit is a
monomorphism in
\[
 0\longrightarrow\mathcal F\longrightarrow\nu_*\nu^*\mathcal F
 \longrightarrow\mathcal Q\longrightarrow0.
\]
Indeed its pullback along the surjection $\nu$ is a split
monomorphism by the adjunction identity, and $\nu^*$ is exact and
conservative. Integral base change at the generic point shows
$\mathcal Q_{\overline\eta}=0$.

Any $p$-torsion sheaf $\mathcal Q$ with zero generic stalk on this
curve has $H^q(R,\mathcal Q)=0$ for $q>n-1$. To verify this
without asserting that $\mathcal Q$ is constructible, write it
as a filtered colimit of constructible $\mathbf F_p$-sheaves
$\mathcal Q_i$. For every $i$, finiteness of its generic module
and the vanishing of the generic colimit give a later $j$ for
which the transition $\mathcal Q_i\to\mathcal Q_j$ is zero
at the generic point. Its image is constructible, with support
in a finite set of closed points. Its cohomology above $n-1$
vanishes. Every class in the cohomological colimit therefore
vanishes in those degrees, proving the assertion.

It follows that if $\cd_p(B)\leq n$, then
$H^q(R,\mathcal F)=0$ for $q>n$. Conversely exact integral
pushforward gives $\cd_p(B)\leq\cd_p(R)$. These comparisons
also hold on every principal localization with the stated
residue-field bounds.

Suppose A.2 holds for $B$. Its good affine open has complement
a finite set of closed points of $\Spec B$. Their images in
$\Spec R$ form a finite closed set, since $\nu$ is integral.
Choose a nonzero $a\in R$ vanishing at these images and, when
$p=2$, also enforcing the preceding sum-of-squares restriction.
Then $\Spec B[1/a]$ lies inside the good open. The upper bound
holds on every smaller open of this normalization, by
Theorem~\ref{A:Dedekind-Milnor-criterion}. For every affine
$U\subset\Spec R[1/a]$, its normalization is the corresponding
open of $\Spec B[1/a]$, so the sheaf comparison proves
$\cd_p(U)\leq n$. Proposition~\ref{A:spread} then proves A.2
for $R$.

Conversely, A.2 for $R$ supplies a principal open with
$\cd_p(R[1/a])\leq n$. Exact integral pushforward gives
$\cd_p(B[1/a])\leq n$. The residue criterion makes this upper
bound hereditary on smaller opens of $\Spec B[1/a]$, and generic
spreading proves A.2 for $B$. The same argument proves the stated
equivalence with the simultaneous residue condition. Finally, if
$R$ is semilocal, the finite-fibre assertion makes $B$ semilocal.
Corollary~\ref{A:semilocal-Dedekind} applies, and the established
normalization comparison finishes the proof.
\end{proof}

\begin{remark}[What the reduction does and does not settle]
This criterion settles every semilocal curve case and makes the obstruction
for every noetherian integral curve explicit. In generic dimension two it is
a simultaneous tame-symbol cokernel, rather than only a Picard
quotient. One finite deletion does not always make these cokernels
vanish: Theorem~\ref{A:A2-persistent-arithmetic-counterexample} gives
an excellent arithmetic counterexample already in generic dimension
one. The Nagata--Gabber example has a different outcome: its
eventual equality follows from the principal-completion theorem
proved below. The literature search used
the combinations ``semilocal'', ``\'etale cohomological dimension'',
``fraction field'', ``Dedekind'', and ``Iwasawa''; it located the
standard norm-residue and arithmetic inputs, without establishing
priority for this formulation.
\end{remark}

\subsection{A principal-completion repair of the Nagata example}

\begin{proposition}[Proved; a quasi-excellent principal quotient suffices]
\label{A:principal-completion-bound}
Let $(B,\mathfrak m)$ be a strictly henselian noetherian local ring
of dimension $d$, let $p$ be invertible in $B$, and let
$t\in\mathfrak m$. If $B/tB$ is quasi-excellent, then
\[
 \cd_p\bigl(\Spec B[1/t]\bigr)\leq d.
\]
No quasi-excellence hypothesis on $B$ is imposed.
\end{proposition}
\begin{proof}
Let $C=\varprojlim_n B/t^nB$ be the $t$-adic completion, and let
$g:\Spec C[1/t]\to\Spec B[1/t]$ be its restriction. The pair
$(B,tB)$ is henselian because $B$ is henselian local and
$tB\subset\mathfrak m$. The Fujiwara--Gabber comparison theorem
in the precise unrestricted form proved by
\citet[Theorem 6.11, p.~50 of version 4]{A_BhattMathew2020}
gives, for every torsion abelian \'etale sheaf $\mathcal F$,
\[
 H^q\bigl(\Spec B[1/t],\mathcal F\bigr)
 \simeq
 H^q\bigl(\Spec C[1/t],g^*\mathcal F\bigr)
 \qquad(q\geq0).
\]
In particular this comparison applies to all $p$-primary torsion
coefficients, without a constructibility restriction.

We verify the hypotheses for the affine vanishing theorem on $C$.
The completion is noetherian. The ideal $tC$ is in its Jacobson
radical, and $C/tC=B/tB$ is local, so $C$ is local with maximal
ideal $\mathfrak m C$. A complete ideal-adic pair is henselian;
combining this with the henselian local quotient $B/tB$ shows that
$C$ is henselian. Its residue field equals that of $B$ and is
separably closed. For every $r$, completion of the finite
$B$-module $B/\mathfrak m^r$ along $t$ leaves it unchanged,
because $t$ is nilpotent on this module. Consequently
$C/(\mathfrak m C)^r=B/\mathfrak m^r$ and the maximal-ideal
completion of $C$ is that of $B$. Thus $\dim C=d$.

The lifting theorem of Nishimura--Nishimura and Gabber,
\citet[Theorem 5.1, p.~7]{A_KuranoShimomoto2019}, says that
a noetherian $I$-adically complete ring is quasi-excellent whenever
its quotient by $I$ is quasi-excellent. Applied to $(C,tC)$,
it makes $C$ quasi-excellent. Gabber's strictly local affine
vanishing theorem
\citep[Expos\'e XV, Corollary 1.2.2]{A_ILO2014}
therefore bounds the cohomological dimension of the affine open
$\Spec C[1/t]$ by $d$. The comparison isomorphism gives the
asserted bound on $\Spec B[1/t]$.
\end{proof}

\begin{theorem}[Proved; A.2 for a nonexcellent class of local surfaces]
\label{A:A2-local-surface}
Let $B$ be a strictly henselian noetherian local domain of dimension
two, and let $p$ be invertible in $B$. Suppose that for some
nonzero $t$ in its maximal ideal the ring $B/tB$ is quasi-excellent.
Then A.2 holds for $\Spec B$ and for every nonempty affine open
subscheme of $\Spec B$.

In particular this applies to every regular strictly henselian
two-dimensional local ring of equicharacteristic zero. For the
Nagata--Gabber ring in Theorem~\ref{A:nagata},
every nonempty affine open of $\Spec B[1/(Ft)]$ has
$p$-cohomological dimension two, even though
$\cd_p(\Spec B[1/F])=3$ for the chosen prime of that example.
\end{theorem}
\begin{proof}
Theorem~\ref{A:strict-field} gives
$\cd_p(\Frac B)=2$. Put $R=B[1/t]$.
Proposition~\ref{A:principal-completion-bound} gives
$\cd_p(R)\leq2$. The ring $R$ has dimension at most one:
the only height-two prime of the two-dimensional local domain $B$
is its maximal ideal, which contains $t$. If $R$ is a field,
the conclusion is immediate. Otherwise its normalization
$R^\nu$ is Dedekind by Krull--Akizuki, and exact integral
pushforward gives $\cd_p(R^\nu)\leq2$.
Theorem~\ref{A:Dedekind-Milnor-criterion} shows that the same
upper bound holds on every open of $\Spec R^\nu$.
The normalization comparison in
Theorem~\ref{A:curve-normalization} transfers the upper bound
to every affine open of $\Spec R$.
Here no extra real-field shrinking is needed: if $p=2$,
strict henselianity and $2\in B^\times$ put a square root of
$-1$ in $B$, and hence in every residue field under consideration.
The lower bound two on every nonempty proper open of $\Spec B$
is again Theorem~\ref{A:strict-field}. Thus every nonempty
affine open contained in $\Spec R$ has $p$-cohomological dimension
exactly two.
Every nonempty open of $\Spec B$ meets $\Spec R$ and contains
a nonempty affine open inside that intersection. This is precisely
the quantifier required in A.2, and it also proves the assertion
for any affine open subscheme of $\Spec B$.

If $B$ is regular of equicharacteristic zero, choose $t$ from a
regular system of parameters. Then $B/tB$ is a characteristic-zero
discrete valuation ring, hence excellent
\citep[Expos\'e I, Proposition 3.1]{A_ILO2014}.
This proves the stated class without assuming that $B$ itself
is excellent. In the explicit Nagata--Gabber construction,
the original element $t$ is such a regular parameter.
The nonempty open $\Spec B[1/(Ft)]$ lies in $\Spec B[1/t]$,
so all its nonempty affine subopens have $p$-cohomological dimension
two.
The dimension-three assertion for $\Spec B[1/F]$ was proved
in Theorem~\ref{A:nagata}.
\end{proof}

\begin{remark}[Scope of the comparison]
The completion in this proof is along the principal ideal $(t)$,
not along the maximal ideal. Maximal-ideal completion compares
punctured spectra, whose general upper bound is three in dimension
two, and by itself would not give the needed affine bound.
The proof does not assert that the original nonexcellent surface
becomes quasi-excellent. It proves the cohomological bound on the
specified principal open, and then justifies its inheritance by
smaller opens using the one-dimensional residue criterion.
\end{remark}

\begin{proposition}[Proved; normal models of a quasi-excellent strictly local field]
\label{A:strict-model-rigidity}
Let $B$ be a quasi-excellent strictly henselian noetherian local
domain, let $p$ be invertible in $B$, and put $K=\Frac B$ and
$d=\dim B$. Every noetherian normal domain $R\subset K$ with
$\Frac R=K$ contains $B$ and satisfies
\[
 \cd_p(\Spec R)\leq d=\cd_p(K).
\]
Consequently A.2 holds for all these normal models, in every
dimension. It also holds for every one-dimensional noetherian
domain with fraction field $K$, without a normality assumption.
\end{proposition}
\begin{proof}
We first prove the containment, which is not assumed in the
statement. For every $u\in B^\times$ and integer $r\geq1$,
strict henselianity lifts a simple root of $X^{p^r}-u$ from the
separably closed residue field. Hence $u$ has a $p^r$th root in
$B^\times$. Any discrete valuation $v:K^\times\to\Z$ must
therefore have $v(u)=0$, since this integer is divisible by
every $p^r$. If $b$ belongs to the maximal ideal of $B$, then
$1+b\in B^\times$. The inequality $v(b)<0$ would force
$v(1+b)=v(b)<0$, a contradiction. Thus every element of $B$
belongs to every discrete valuation ring of $K$.

For every height-one prime $\mathfrak q$ of $R$, the ring
$R_{\mathfrak q}$ is a discrete valuation ring. The intersection
formula for a noetherian normal domain
\citep[Tag 031T, Lemma 10.157.6(2)]{A_Stacks} yields
\[
 B\subset\bigcap_{\operatorname{ht}\mathfrak q=1}
 R_{\mathfrak q}=R.
\]
If $R$ is a field, the containment follows from $R=K$.

Write $R$ as the filtered union of its finite-type $B$-subalgebras
$R_i$. Each $R_i$ is a domain, the morphism
$\Spec R_i\to\Spec B$ is affine, finite type and dominant,
and its generic relative dimension is zero, since
$\Frac R_i=K$. The affine bound of
Theorem~\ref{A:relative-affine} gives $\cd_p(R_i)\leq d$.
Descent of constructible sheaves and continuity of their
cohomology, followed by passage to all torsion coefficients as
in Lemma~\ref{A:A2-approximation}, give $\cd_p(R)\leq d$.
The equality for $K$ is Theorem~\ref{A:strict-field}.

Every affine open of $\Spec R$ again has a noetherian normal
coordinate ring with fraction field $K$, so the same upper bound
applies. The generic lower witness from
Proposition~\ref{A:spread} supplies A.2. Finally, for any
one-dimensional noetherian domain with this fraction field,
its normalization is a noetherian Dedekind domain with the same
field; the normal case and Theorem~\ref{A:curve-normalization}
give A.2 without assuming normality.
\end{proof}

\subsection{Arithmetic approximation in finite transcendence degree}

\begin{lemma}[Proved; the affine bound over an arithmetic curve]
\label{A:arithmetic-affine-relative}
Let $F$ be a number field, let $p$ be a prime, and let
$S=\Spec\mathcal O_F[1/p]$. Suppose that $p$ is odd or that
$F$ is totally imaginary. Let $f:T\to S$ be an affine, flat
morphism of finite type whose generic fibre has dimension at
most $e$. Then
\[
 \cd_p(T)\leq e+2.
\]
The scheme $T$ need not be integral or normal.
\end{lemma}
\begin{proof}
The base has $p$-cohomological dimension at most two, by
Artin--Verdier duality in the form used in
Proposition~\ref{A:arithmetic-test}. For odd $p$ the real-place
terms have no positive $p$-cohomology; thus the same upper bound
holds without the totally imaginary hypothesis. Every closed
point of $S$ has a finite residue field and hence
$p$-cohomological dimension one.

Give $S$ its dimension function $\delta_S$ with value one at
the generic point and zero at every closed point. Flatness
excludes vertical irreducible components of $T$. The dimension
formula for finite-type schemes over the excellent regular
one-dimensional scheme $S$ shows that all fibres have dimension
at most $e$. Therefore the induced function
\[
 (f^*\delta_S)(t)
 =\delta_S(f(t))+\operatorname{trdeg}_{\kappa(f(t))}\kappa(t)
\]
has supremum at most $e+1$. This argument applies to every
irreducible component and makes no integrality assumption on $T$.

Let $\mathcal F$ be a constructible sheaf of
$\mathbf F_p$-vector spaces on $T$. Gabber's affine theorem
\citep[Expos\'e XV, Theorem 1.1.2 and its reformulation 1.2.1,
p.~219]{A_ILO2014} gives
\[
 \delta_S\bigl(R^b f_*\mathcal F\bigr)\leq e+1-b.
\]
Thus $R^b f_*\mathcal F=0$ for $b>e+1$, and
$\mathcal Q=R^{e+1}f_*\mathcal F$ has zero geometric generic
stalk. The latter conclusion can also be read from the ordinary
affine theorem on the geometric generic fibre.

We check the cohomology of this top direct image using the
generic-zero sheaf argument in the proof of
Theorem~\ref{A:curve-normalization}. Write
$\mathcal Q=\varinjlim_i\mathcal Q_i$ with each
$\mathcal Q_i$ constructible. For each $i$, finiteness of its
generic stalk and vanishing of the generic colimit give a later
$j$ for which $\mathcal Q_i\to\mathcal Q_j$ is zero at the
generic point. The image of this transition is constructible
and supported on finitely many closed points. Exact closed
pushforward and the dimension-one bound for finite fields give
vanishing above degree one for this image. The transition on
cohomology factors through that image, so every class vanishes
at a later stage. Cohomological continuity gives
$H^a(S,\mathcal Q)=0$ for $a>1$, whether or not constructibility
of $\mathcal Q$ is invoked.

In the Leray spectral sequence
\[
 H^a\bigl(S,R^b f_*\mathcal F\bigr)
 \Longrightarrow H^{a+b}(T,\mathcal F),
\]
the rows $b\leq e$ vanish for $a>2$, the row $b=e+1$ vanishes
for $a>1$, and higher rows vanish identically. Hence the total
cohomology vanishes in degrees greater than $e+2$.
Coefficient d\'evissage and constructible approximation on the
noetherian scheme $T$ prove the bound for all $p$-primary torsion
sheaves.
\end{proof}

\begin{proposition}[Proved; arithmetic approximation beyond finite type]
\label{A:arithmetic-transcendence}
Let $R$ be a characteristic-zero domain, let $p$ be invertible
in $R$, and suppose
$d=\operatorname{trdeg}_{\Q}\operatorname{Frac}R<\infty$.
No noetherian hypothesis is imposed. If $p$ is odd, then
\[
 \cd_p(\Spec R)\leq d+2.
\]
If $p=2$ and $K=\operatorname{Frac}R$ is not orderable, there
is a nonzero $a\in R$ such that
\[
 \cd_2\bigl(\Spec R[1/a]\bigr)\leq d+2.
\]
The same bounds apply to every affine open of the indicated scheme.
\end{proposition}
\begin{proof}
For odd $p$, write $R$ as the filtered union of its finite-type
$\Z[1/p]$-subalgebras $R_i$. These are domains and are flat over
$\Z[1/p]$, since they are torsion free as modules over a
Dedekind domain. Their generic fibres have dimensions
$e_i=\operatorname{trdeg}_{\Q}\operatorname{Frac}(R_i)\leq d$.
Lemma~\ref{A:arithmetic-affine-relative} gives
$\cd_p(R_i)\leq d+2$ uniformly. Descent of constructible
coefficients, continuity of cohomology, and the passage to all
torsion sheaves in Lemma~\ref{A:A2-approximation} give the same
bound on $R$.

Now let $p=2$. Write $-1=\sum_{j=1}^m x_j^2$ in $K$ and
choose a common nonzero denominator $a$ so that each
$x_j\in R[1/a]$. Approximate $R[1/a]$ by its finite-type
$\Z[1/2]$-subalgebras $A_i$ containing all $x_j$.
Every residue field of $A_i$ is nonorderable, because the displayed
identity survives in it. The fraction field of $A_i$ is a
finitely generated nonorderable characteristic-zero field. It has
finite $2$-cohomological dimension, by the field theorem
\citep[Expos\'e X, Theorem 2.1]{A_SGA4}. Consequently
Corollary~\ref{A:generic-plus-dimension} gives
$\cd_2(A_i)<\infty$.

Use the quadratic finite \emph{\'etale cover}
\[
 A_i'=A_i[T]/(T^2+1),\qquad
 \Spec A_i'\longrightarrow\Spec A_i.
\]
We do not assume that $A_i'$ is a domain. As a scheme over
$\Spec\Z[i,1/2]$, it is the base change of the flat
$\Spec A_i\to\Spec\Z[1/2]$. It is therefore affine, flat
and of finite type, and its generic fibre has dimension at most
$d$. Lemma~\ref{A:arithmetic-affine-relative}, with the totally
imaginary number field $\Q(i)$, gives $\cd_2(A_i')\leq d+2$.
The cover is a torsor under the constant group of order two, so
Lemma~\ref{A:quadratic-descent}, using the finite bound already
proved for $A_i$, gives $\cd_2(A_i)\leq d+2$.
The same continuity argument passes this uniform estimate to
$R[1/a]$.

Finally, the coordinate ring of a nonempty affine open has the
same fraction field. For odd $p$ the first argument applies to it
directly. In the second case it contains the elements $x_j$, so
the approximation by subalgebras containing the one fixed
sum-of-squares identity applies without a further localization.
This proves the last assertion.
\end{proof}

\begin{corollary}[Proved; arbitrary noetherian models of maximal arithmetic dimension]
\label{A:A2-arithmetic-maximal}
Let $R$ be a noetherian characteristic-zero domain, let $p$ be
invertible in $R$, and put $K=\operatorname{Frac}R$. If
\[
 \cd_p(K)=\operatorname{trdeg}_{\Q}K+2<\infty,
\]
then A.2 holds for $\Spec R$. In particular it holds for
\emph{every} noetherian model whose fraction field is finitely
generated over $\Q$ and whose $p$-cohomological dimension is
finite. Excellence, normality and finite type of the model
are unnecessary.
\end{corollary}
\begin{proof}
For $p=2$ finite cohomological dimension excludes an ordering of
$K$. Proposition~\ref{A:arithmetic-transcendence} therefore
supplies an open $V$ whose every affine open has cohomological
dimension at most $n=\cd_p(K)$; for odd $p$ we may take
$V=\Spec R$. Given a nonempty open $W$, apply
Proposition~\ref{A:spread} inside $W\cap V$. Its lower witness
survives on every smaller nonempty affine open, giving equality
there. This is A.2 with its full quantifier over smaller opens.

For a field finitely generated over $\Q$, the field theorem
\citep[Expos\'e X, Theorem 2.1]{A_SGA4} gives the displayed
equality for odd $p$. For $p=2$, apply it after adjoining $i$;
finite extension preserves finite $2$-cohomological dimension
for the nonorderable field $K$. Thus the last assertion has exactly
the stated hypothesis.
\end{proof}

\begin{remark}[The remaining arithmetic range]
The bound is obtained from the opened affine and arithmetic
duality theorems; we make no priority claim. We searched the
phrases ``affine cohomological dimension arithmetic schemes'',
``cohomological dimension finite type over integers'', and
``dimension cohomologique affine arithm\'etique'', and checked
SGA~4, Expos\'es X and XIV, and the precise affine statement in
\citet[Expos\'e XV, Theorem 1.1.2]{A_ILO2014}.
This reduction leaves, among finite-transcendence-degree
characteristic-zero fields, those for which the cohomological
dimension drops below $\operatorname{trdeg}_{\Q}K+2$.
The cyclotomic arithmetic instance remains in that range and
is not settled by this approximation argument.
\end{remark}

\subsection{A restriction on arithmetic towers with noetherian integral closure}

\begin{proposition}[Proved; the Galois tower is virtually procyclic]
\label{A:A2-Galois-tower-procyclic}
Let $F$ be a number field, let $S$ be a finite set of finite primes
of $F$, and let $K/F$ be an algebraic Galois extension.  Let $R$ be
the integral closure of $\mathcal O_F[S^{-1}]$ in $K$.  If $R$ is
noetherian, then $\operatorname{Gal}(K/F)$ has an open procyclic
subgroup.  In particular, if $K/F$ is an infinite pro-$p$ extension,
there is a finite intermediate field $F'/F$ such that $K/F'$ is a
$\mathbf Z_p$-extension.
\end{proposition}

\begin{proof}
Write $A=\mathcal O_F[S^{-1}]$.  The ring $R$ is a normal domain,
integral over the Dedekind domain $A$, with fraction field $K$.
Its nonzero prime ideals are maximal.  Thus, under the stated
noetherianity hypothesis, $R$ is Dedekind.

Fix a nonzero prime $\mathfrak q$ of $A$.  Choose
$0\ne a\in\mathfrak q$.  Every prime of $R$ over $\mathfrak q$
contains $a$, so the factorization of $aR$ in the Dedekind domain
$R$ shows that there are only finitely many such primes.  Choose
one, $Q$, and let $D\subseteq G=\operatorname{Gal}(K/F)$ be its
decomposition subgroup.  The Galois action on the primes over
$\mathfrak q$ is transitive, by the finite-extension statement
and compactness.  Consequently $D$ is open in $G$.

We also record explicitly the ramification consequence of
noetherianity.  The local ring $R_Q$ is a discrete valuation ring;
normalize its valuation $v_Q$ to have value group $\mathbf Z$.
If $\pi\in F$ is a uniformizer at $\mathfrak q$, put
$e_\infty=v_Q(\pi)>0$.  For a finite intermediate field
$F\subseteq E\subseteq K$ and the induced prime $Q_E$, the
restriction of $v_Q$ to $E^\times$ is a positive integral
multiple of the normalized valuation at $Q_E$.  It follows that
\[
                 e(Q_E/\mathfrak q)\mid e_\infty.
\]
In particular, all finite-level ramification indices along $Q$
are bounded by the same integer.

For finite Galois intermediate fields, the inertia group has
order equal to that ramification index, and the quotient of the
decomposition group by inertia is the Galois group of the finite
residue-field extension.  We use the local decomposition and
inertia statements in \cite[Theorem~7.58 and
Propositions~8.10--8.11, pp.~131, 139--140]{C_MilneANT2020}.
Passing to inverse limits therefore gives an exact sequence
\[
  1\longrightarrow I\longrightarrow D\longrightarrow
  \operatorname{Gal}(\kappa(Q)/\kappa(\mathfrak q))
  \longrightarrow1,
\]
where $I$ is finite: all its finite quotient orders are bounded
by $e_\infty$.  The group on the right is procyclic because
$\kappa(\mathfrak q)$ is finite.

Choose a topological generator $\overline\sigma$ of $D/I$ and
a lift $\sigma\in D$.  The closed subgroup
$H=\overline{\langle\sigma\rangle}$ is procyclic.  Its image
in $D/I$ is compact and contains the dense cyclic subgroup
generated by $\overline\sigma$, so that image is all of $D/I$.
Hence $D=HI$ and $[D:H]\le |I|$.  A closed finite-index
subgroup of a profinite group is open, so $H$ is open in $G$.
If $G$ is infinite pro-$p$, then $H$ is an infinite procyclic
pro-$p$ group, hence isomorphic to $\mathbf Z_p$.  Its fixed
field is the required finite extension $F'$.
\end{proof}

\begin{remark}[Exact scope of the obstruction]
We did not assume that ramification is confined to $S$, or even
that there are only finitely many ramified primes.  One retained
prime already gives the conclusion.  The proposition concerns the
\emph{full integral closure} in a Galois tower.  A non-Galois
tower, or an overring obtained by discarding a non-Galois-invariant
collection of primes, is outside its conclusion.  Thus a
Golod--Shafarevich Galois tower with non-virtually-procyclic group
cannot directly furnish the sought noetherian Dedekind ring by
taking its full ring of $S$-integers.
\end{remark}

\begin{lemma}[Proved; the exact persistence requirement]
\label{A:A2-class-persistence}
Let $F_0\subset F_1\subset\cdots$ be an algebraic-number-field
tower, let $S_n$ be the primes above a fixed finite set $S_0$,
and put $A_n=\mathcal O_{F_n}[S_n^{-1}]$ and
$A_\infty=\varinjlim_n A_n$.  Then
\[
 \operatorname{Pic}(A_\infty)/p
   =\varinjlim_n\bigl(\operatorname{Pic}(A_n)/p\bigr),
\]
where the transition maps extend ideals.  A class
$c\in\operatorname{Pic}(A_n)$ survives nontrivially in this
quotient if and only if its extension to
$\operatorname{Pic}(A_m)$ belongs to no subgroup
$p\operatorname{Pic}(A_m)$, for every $m\ge n$.
More generally, a finite list is linearly independent in the
limit if and only if no nonzero $\mathbf F_p$-linear combination
of it becomes $p$-divisible at any later stage.
\end{lemma}

\begin{proof}
Invertible modules and their isomorphisms descend through a
filtered colimit of rings: we descend a finite projective
presentation, its rank-one condition, and an inverse module with
its evaluation isomorphism.  Therefore
$\operatorname{Pic}(A_\infty)=\varinjlim_n\operatorname{Pic}(A_n)$.
Taking the cokernel of multiplication by $p$ commutes with the
filtered colimit.  The stated survival and independence tests
are the defining zero test in a filtered colimit of vector spaces.
\end{proof}

If $A_\infty$ is Dedekind and this vector space is infinite
dimensional, deleting any finite set of closed points leaves it
nonzero: localization quotients its Picard group by the subgroup
generated by those finitely many prime classes.  By contrast,
unbounded finite-stage class-group ranks, or injectivity before
reduction modulo $p$, do not establish the persistence condition.


\paragraph{The non-Galois construction is completed below.}
Lemma~\ref{A:A2-quadratic-step} combines unramified inertness at
any prescribed finite set of primes with injectivity of ideal
extension modulo two and growth of its dimension. The resulting
non-Galois tower in Theorem~\ref{A:A2-persistent-arithmetic-counterexample}
satisfies both the noetherianity and class-persistence requirements.
The virtually procyclic restriction above explains why the full
integral closure in an arbitrary Galois tower was an inadequate route.

\subsection{Dedekind overrings and all models of certain strictly local fields}

\begin{lemma}[Proved; the generic upper bound passes to Dedekind overrings]
\label{A:Dedekind-overring-cd}
Let $D$ be a Dedekind domain, let $p\in D^\times$ be prime,
and let $K=\Frac D$ have finite $p$-cohomological dimension $n$.
Suppose $\cd_p(\Spec D)\leq n$. Then every ring
$D\subset E\subset K$ satisfies
\[
 \cd_p(\Spec E)\leq n.
\]
Every such $E$ is a Dedekind domain or a field. It need not be a
localization of $D$ at a multiplicative subset.
\end{lemma}
\begin{proof}
If $D$ is a field there is nothing to prove. Otherwise $n\geq1$.
Krull--Akizuki, in the overring form of
\citet[Tag 00PG, Lemma 10.119.12]{A_Stacks}, makes $E$ noetherian.
We recall explicitly the prime correspondence that we need.
For a nonzero prime $\mathfrak q\subset E$, its contraction
$\mathfrak p=\mathfrak q\cap D$ is nonzero: writing a nonzero
$x\in\mathfrak q$ as $a/b$ with nonzero $a,b\in D$ gives
$a=bx\in\mathfrak p$. The inclusion
$D_{\mathfrak p}\subset E_{\mathfrak q}\subset K$ is an inclusion
of local overrings of the discrete valuation ring $D_{\mathfrak p}$,
and its uniformizer remains a nonunit. The only proper overring
of a discrete valuation ring in its fraction field is the field.
Consequently
\[
 E_{\mathfrak q}=D_{\mathfrak p}.
\]
This proves that $E$ is Dedekind when it is not a field. It also
shows that contraction is injective on nonzero primes, with
unchanged local discrete valuation rings and residue fields:
$\mathfrak q$ is recovered as the inverse image of the maximal
ideal of this common local ring. This is the overring correspondence
also stated in \citet[Lemma 1, p.~216]{A_Clark2009}.

Fix a finite separable extension $L/K$ containing $\mu_p$.
Let $D_L$ and $E_L$ be the integral closures of $D$ and $E$ in $L$.
They are Dedekind or fields and are finite over their respective
bases, by the trace argument in
Lemma~\ref{A:Dedekind-Milnor-obstruction}.
Since every element integral over $D$ is integral over $E$,
$D_L\subset E_L\subset L$. The preceding local-ring argument
identifies the nonzero primes of $E_L$ with a subset of those of
$D_L$, preserving the associated normalized discrete valuations.
The residue map for $E_L$ is therefore the composite
\[
 K_n^M(L)/p\longrightarrow
 \bigoplus_{\mathfrak p\in(\Spec D_L)^{(1)}}
 K_{n-1}^M(\kappa(\mathfrak p))/p
 \longrightarrow
 \bigoplus_{\mathfrak q\in(\Spec E_L)^{(1)}}
 K_{n-1}^M(\kappa(\mathfrak q))/p,
\]
where the second arrow is the projection to that subset.
The first arrow is surjective by
Theorem~\ref{A:Dedekind-Milnor-criterion} and the hypothesis on $D$;
hence the composite is surjective. The same criterion applied to
$E$ gives the asserted bound. If $E=K$, the conclusion is immediate.
No finite-type hypothesis on $D\to E$ has entered the argument.
\end{proof}

\begin{lemma}[Proved; strictly local rings lie in every noetherian model normalization]
\label{A:strict-normalization-rigidity}
Let $B$ be a strictly henselian noetherian local domain with a prime
$p$ invertible in $B$, and put $K=\Frac B$. If $R\subset K$ is any
noetherian domain with fraction field $K$, its integral closure
$\overline R$ in $K$ contains $B$. The normalization need not be
finite or noetherian.
\end{lemma}
\begin{proof}
Every $u\in B^\times$ has a $p^r$th root in $B^\times$ for each
$r\geq1$: the polynomial $X^{p^r}-u$ has a simple root over the
separably closed residue field and strict henselianity lifts it.
Thus every discrete valuation $v:K^\times\to\Z$ vanishes on
$B^\times$, since $v(u)$ is divisible by every $p^r$.
For $b$ in the maximal ideal of $B$, the element $1+b$ is a unit.
If $v(b)<0$, the valuation inequality with unequal values gives
$v(1+b)=v(b)<0$, a contradiction. Therefore $B$ lies in every
discrete valuation ring of $K$.

The Mori--Nagata theorem
\citep[Theorem 4.10.5, pp.~90--92]{A_SwansonHuneke2006}
states that the integral closure of a noetherian domain in its fraction field is a Krull ring. In particular,
$\overline R$ is the intersection in $K$ of its height-one
discrete valuation local rings. The preceding paragraph puts $B$
in each of these rings, and hence in $\overline R$. If
$\overline R=K$, the same containment is immediate. We have not
replaced a possibly infinite normalization by a finite one. The same global formulation is recorded in 
\citet[Chapitre 0, Remarque 23.2.8(iii), p.~220]{A_EGAIV1}.
\end{proof}

\begin{theorem}[Proved; A.2 for every model of a specified strictly local surface field]
\label{A:A2-all-surface-models}
Let $B$ be a strictly henselian noetherian local domain of dimension
two, let $p$ be invertible in $B$, and suppose that some nonzero
$t$ in its maximal ideal has quasi-excellent quotient $B/tB$.
Put $K=\Frac B$. Then A.2 holds for every noetherian domain
$R\subset K$ with $\Frac R=K$ and $p\in R^\times$.
The domain $R$ need not be normal, excellent, local, or of finite
type over a subring of $B$.

In particular, this covers every noetherian model of the fraction
field of every regular strictly henselian local surface of
equicharacteristic zero, including the Nagata--Gabber example.
\end{theorem}
\begin{proof}
We know $\cd_p(K)=2$ by Theorem~\ref{A:strict-field}.
The ring $B[1/t]$ has dimension at most one, and
Proposition~\ref{A:principal-completion-bound} gives
$\cd_p(B[1/t])\leq2$. Let $D$ be its integral closure in $K$.
Krull--Akizuki makes $D$ Dedekind or a field. Exact integral
pushforward on the small \'etale site gives
\[
 \cd_p(\Spec D)\leq2.
\]
Indeed, for the integral map $\nu:\Spec D\to\Spec B[1/t]$
and every $p$-primary torsion sheaf $\mathcal E$ on $\Spec D$,
$H^q(\Spec D,\mathcal E)=
H^q(\Spec B[1/t],\nu_*\mathcal E)$.

Let $\overline R$ be the integral closure of $R$ in $K$.
Lemma~\ref{A:strict-normalization-rigidity} gives
$B\subset\overline R$. Since $t\in K^\times$, choose
$0\neq a\in R$ so that both $t$ and $t^{-1}$ belong to $R[1/a]$.
For example, if $t=r/s$ with nonzero $r,s\in R$, we may initially
take $a=rs$. When $p=2$, strict henselianity puts an element
$i$ with $i^2=-1$ in $B$; enlarge $a$ by the denominator of an
expression for $i$ in $K=\Frac R$. Thus all residue fields of
$R[1/a]$ are nonorderable in this case.

Write $R_0=R[1/a]$ and $E=\overline R[1/a]$.
Integral closure commutes with localization, so $E$ is the
integral closure of $R_0$. It is integrally closed in $K$ and
contains $B[1/t]$, and therefore contains $D$.
Lemma~\ref{A:Dedekind-overring-cd} gives that $E$ is Dedekind
or a field and that $\cd_p(E)\leq2$.
The integral surjection $\Spec E\to\Spec R_0$ preserves
Krull dimension, so $\dim R_0\leq1$.

More generally, let $W\subset\Spec R_0$ be a nonempty affine
open and put $S=\Gamma(W,\mathcal O_W)\subset K$.
The normal ring $S^\nu$, the integral closure of $S$ in $K$,
contains $E$: every element of $E$ is integral over $R_0\subset S$.
Thus it is again an overring of $D$ and has cohomological
dimension at most two. If $S$ is a field then $S=K$ and the same
bound holds. Otherwise $S$ is a noetherian one-dimensional domain.
The sheaf comparison proved in
Theorem~\ref{A:curve-normalization} gives $\cd_p(S)\leq2$:
the unit into integral-normalization pushforward has a generically
zero cokernel whose cohomology vanishes above degree one.
The real-residue exception has already been excluded when $p=2$.
Hence the upper bound two holds on every such affine $W$.

Finally Proposition~\ref{A:spread} supplies a nonempty open on
which every nonempty affine subopen has cohomological dimension
at least $\cd_p(K)=2$. Intersect it with $\Spec R_0$.
Every nonempty open of $\Spec R$ contains a nonempty affine open
inside this intersection, and every further nonempty affine
subopen has cohomological dimension exactly two. This proves the
full quantifier in A.2.

For the last assertion, a regular parameter $t$ in a regular
strictly henselian local surface of equicharacteristic zero has
characteristic-zero discrete valuation quotient, hence excellent,
as checked in Theorem~\ref{A:A2-local-surface}.
\end{proof}

\begin{remark}[Search and scope]
We checked the full Mori--Nagata proof in Huneke--Swanson,
the Stacks Project's precise Krull--Akizuki overring statement,
and Clark's explicit Dedekind prime correspondence. The
cohomological comparison uses the residue criterion already proved
above. We make no priority claim for the resulting model theorem.
It removes all noetherian models of these specified surface fields
from the A.2 gate; it does not settle arbitrary one-dimensional
arithmetic fields or arbitrary strictly local fields in higher
dimension.
\end{remark}

\subsection{Topological normalization and semilocal models}

\begin{lemma}[Known; finite topological normalization, with the semilocal derivation]
\label{A:semilocal-topnormalization}
Let $R$ be a semilocal noetherian domain with fraction field $K$,
and let $\overline R$ be its integral closure in $K$.
There is a finite birational $R$-subalgebra
$R\subset C\subset\overline R$ such that
\[
 \Spec\overline R\longrightarrow\Spec C
\]
is a universal homeomorphism. Consequently there is $0\neq a\in R$
such that the normalization morphism of $\Spec R[1/a]$ is a
universal homeomorphism. Neither normalization is assumed finite
or noetherian.
\end{lemma}
\begin{proof}
Let $\mathfrak m_1,\ldots,\mathfrak m_r$ be the maximal ideals of
$R$ and set
\[
 T=\prod_{i=1}^r\widehat{R_{\mathfrak m_i}},
\]
where each completion is along its maximal ideal. The map $R\to T$
is flat, since localization and noetherian completion are flat and
the product is finite. It is faithfully flat: for a proper ideal
$I\subset R$, choose $\mathfrak m_i\supset I$; the corresponding
factor of $T/IT$ has a nonzero residue field. The ring $T$ is
noetherian. Each of its minimal-prime quotients is a complete
noetherian local domain, and hence is Japanese by
\citet[Chapitre 0, Th\'eor\`eme 23.1.5, pp.~215--216]{A_EGAIV1}.

Apply \citet[Chapitre 0, Proposition 23.2.4(ii), p.~218]{A_EGAIV1}
to the faithfully flat map $R\to T$ and the trivial finite
extension $K/K$. Its directed family consists of the finite
$R$-subalgebras of $\overline R$; each already has fraction field
$K$ because it contains $R$. The proposition gives such an
algebra $C$ for which $\Spec\overline R\to\Spec C$ is radicial.
This map is also integral and surjective, so it is a universal
homeomorphism. For a local ring this is exactly the consequence
stated in Corollaire 23.2.5, pp.~218--219, of the same source.

Choose finite $R$-algebra generators $c_1,\ldots,c_s$ of $C$ and
write them as fractions in $K=\Frac R$. The product $a\neq0$ of
their denominators gives $C[1/a]=R[1/a]$. Integral closure
commutes with localization, and universal homeomorphisms are
preserved by base change. Therefore
\[
 \Spec\overline R[1/a]\longrightarrow\Spec R[1/a]
\]
is the asserted normalization and universal homeomorphism.
\end{proof}

\begin{lemma}[Proved; the strictly local approximation bound for every overring]
\label{A:strict-overring-bound}
Let $B$ be a quasi-excellent strictly henselian noetherian local
domain, let $p$ be an invertible prime, and write
$K=\Frac B$ and $d=\dim B$. Every intermediate ring
$B\subset T\subset K$ satisfies
\[
 \cd_p(\Spec T)\leq d.
\]
The ring $T$ need not be noetherian or integrally closed.
\end{lemma}
\begin{proof}
Write $T$ as the filtered union of its finite-type $B$-subalgebras
$T_i$. Each $T_i$ is a domain containing $B$ and contained in $K$,
so $\Frac T_i=K$. Thus $\Spec T_i\to\Spec B$ is affine,
finite type and dominant with generic relative dimension zero.
The strictly local affine theorem, in the form proved and checked
in Theorem~\ref{A:relative-affine}, gives
$\cd_p(\Spec T_i)\leq d$.

The inverse system of affine schemes $\Spec T_i$ has affine
transition maps and limit $\Spec T$. A constructible torsion
\'etale sheaf descends to a finite stage, and its cohomology is
the filtered colimit of the cohomology at subsequent stages.
Consequently it has zero cohomology in degrees greater than $d$.
The coherent-site constructible approximation and coefficient
d\'evissage used in Lemma~\ref{A:A2-approximation} then give the
same vanishing for every $p$-primary torsion sheaf on $\Spec T$.
This is precisely the argument of
Proposition~\ref{A:strict-model-rigidity}, with its unnecessary
noetherianity and normality hypotheses on the intermediate ring
removed; the base $B$ still has all the hypotheses required by
the affine theorem.
\end{proof}

\begin{theorem}[Proved under stated hypotheses; A.2 for topologically normal and semilocal models]
\label{A:A2-topnormal-models}
Let $B,p,K,d$ be as in Lemma~\ref{A:strict-overring-bound}.
Let $R\subset K$ be a noetherian domain with $\Frac R=K$ and
$p\in R^\times$. Suppose that $\Spec R$ has a nonempty
geometrically unibranch open subscheme. Then A.2 holds for
$\Spec R$. In particular, A.2 holds for every semilocal
noetherian model $R$ of this field, without a normality assumption
and in every dimension.
\end{theorem}
\begin{proof}
Theorem~\ref{A:strict-field} gives $\cd_p(K)=d$.
Choose a nonzero $a\in R$ such that $\Spec R[1/a]$ lies in the
given geometrically unibranch open. For an irreducible
geometrically unibranch scheme the normalization is a universal
homeomorphism, without any finiteness hypothesis on the
normalization; this is
\citet[Tag 0GIQ, Lemma 29.55.12]{A_Stacks}.

Let $W\subset\Spec R[1/a]$ be any nonempty affine open and put
$S=\Gamma(W,\mathcal O_W)\subset K$. Its normalization
$S^\nu\subset K$ is the normalization restricted to $W$, since
normalization commutes with restriction to open subschemes.
Consequently $\Spec S^\nu\to\Spec S$ is a universal
homeomorphism. The Mori--Nagata containment of
Lemma~\ref{A:strict-normalization-rigidity}, applied to the
noetherian model $S$, gives
\[
 B\subset S^\nu\subset K.
\]
By Lemma~\ref{A:strict-overring-bound},
$\cd_p(\Spec S^\nu)\leq d$. Topological invariance of the
small \'etale site
\citep[Tag 04DZ, Theorem 59.45.2]{A_Stacks}
is an equivalence of sites for this integral universal
homeomorphism. It therefore compares all torsion coefficients
and gives
\[
 \cd_p(W)=\cd_p(\Spec S^\nu)\leq d.
\]
No replacement of an infinite normalization by a finite one has
been made in this comparison.

Proposition~\ref{A:spread} supplies a nonempty affine open on
which a nonzero generic degree-$d$ class spreads and remains
nonzero on every smaller nonempty affine open. Intersect with
$\Spec R[1/a]$ and with any prescribed nonempty open of
$\Spec R$, and choose a nonempty affine open $U$ in this
intersection. Every nonempty affine $U'\subset U$ then has
cohomological dimension at least $d$ by that generic class and at
most $d$ by the preceding paragraph. This proves all the
quantifiers of A.2. If $d=0$, the degree-zero generic class gives
the same conclusion.

For semilocal $R$, Lemma~\ref{A:semilocal-topnormalization}
provides a principal open whose normalization is a universal
homeomorphism. The preceding proof applies directly to that
open. Equivalently, the local-ring definition of geometric
unibranchness shows that this principal open has that property.
\end{proof}

\begin{remark}[Exact additional range]
This extends the normal-model result to geometrically unibranch
dense opens, and in particular to nonnormal semilocal models in
dimensions greater than two. It does not make the analogous
claim for every noetherian model. The finite-product argument
in Lemma~\ref{A:semilocal-topnormalization} uses semilocality:
an infinite product of local completions need not be noetherian,
so the cited EGA proposition cannot be applied to it. The
scope of these positive results is recorded in
Section~\ref{A:unresolved}. The unrestricted A.2 assertion is refuted
by the separate arithmetic construction below.
\end{remark}

\subsection{An arithmetic counterexample to eventual equality in A.2}

\begin{lemma}[Proved; the elementary two-group calculation]
\label{A:A2-wreath-calculation}
Let $V$ be a nonzero finite-dimensional vector space over $\mathbf F_2$.
Put $G=(V\oplus V)\rtimes\langle\tau\rangle$, where $\tau$ interchanges
the two factors, put $H=V\oplus V$, and define
\[
 \Sigma:G\longrightarrow V,\qquad (x,y;\epsilon)\longmapsto x+y.
\]
Then $A=\ker\Sigma$ is an elementary abelian two-group. Writing its
elements as $(b,u)\leftrightarrow(u,u;b)$ and using the section
$s(v)=(v,0;0)$ of $\Sigma$, the conjugation action is
\[
                 v\cdot(b,u)=(b,u+bv).
\]
There is a surjective $V$-module homomorphism
\[
 \phi:\mathbf F_2[V]\longrightarrow A,
 \qquad e_g\longmapsto(1,g),
\]
and hence a surjection $\mathbf F_2[V]\rtimes V\to G$ inducing the
identity on the quotient $V$. Each $\phi(e_g)$ is a conjugate of
$\tau$, and its cyclic subgroup intersects $H$ trivially.
The transfer from $G$ to its index-two abelian subgroup $H$ is
\[
             \operatorname{Ver}_{G,H}(x,y;\epsilon)
                    =(x+y,x+y).
\]
\end{lemma}
\begin{proof}
Multiplication on $A$ is componentwise addition because the two entries
of $(u,u)$ agree. Conjugation by $s(v)$ sends $(u,u;b)$ to
$(u+bv,u+bv;b)$, proving the action formula. Translation by $v$ sends
$e_g$ to $e_{g+v}$, which proves equivariance. The image contains
$(1,0)$ and $(1,g)+(1,0)=(0,g)$ for every $g$, so it is all of $A$.
The conjugates of $\tau$ by $(g,0;0)$ are $(g,g;1)$, precisely the
stated images. Their nonidentity elements have nonzero image in
$G/H\simeq C_2$.

For the transfer, use coset representatives $1,\tau$. An element
$(x,y;0)$ transfers to the product with its $\tau$-conjugate,
namely $(x+y,x+y)$. An element $(x,y;1)$ exchanges the two cosets;
its transfer is its square, again $(x+y,x+y)$. These calculations
also prove that in $A\rtimes V$ an element $(b,u;v)$ has square
$(0,bv;0)$.
\end{proof}

\begin{lemma}[Proved; a quadratic step preserving ideal classes modulo two]
\label{A:A2-quadratic-step}
Let $F$ be a totally imaginary number field, and let $S$ be a finite
set of its finite primes containing every prime above $2$. Write
\[
 C_F=\operatorname{Pic}(\mathcal O_F[1/2])/2,
 \qquad r=\dim_{\mathbf F_2} C_F>0.
\]
There is a quadratic extension $L/F$ with the following properties:
\begin{enumerate}
 \item Every prime in $S$ is unramified and inert in $L/F$.
 \item Extension of invertible ideals induces an injection
       $C_F\hookrightarrow C_L$.
 \item $\dim_{\mathbf F_2}C_L\ge 2r$.
\end{enumerate}
More precisely, we can construct a Galois extension $N/F$ containing
$L$ such that
\[
 \operatorname{Gal}(N/F)\simeq (V\oplus V)\rtimes C_2,
 \qquad \operatorname{Gal}(N/L)=V\oplus V,
 \qquad V\simeq C_F,
\]
and $N/L$ is unramified at every finite prime and splits completely
at every prime above $2$.
\end{lemma}
\begin{proof}
Let $M/F$ be the maximal elementary abelian two-extension unramified
at all finite primes and split at the primes above $2$. Global class
field theory identifies $V=\operatorname{Gal}(M/F)$ with $C_F$ by
the Artin map. We use the ideal-class existence and reciprocity
theorems in \cite[Chapter V, Theorems 3.5--3.6, p.~158]{C_MilneCFT2020},
and take the quotient in which the classes of the dyadic primes
vanish. Thus $M/F$ is everywhere unramified and $2$-split.

For each $\mathfrak q\in S$, choose one prime
$\mathfrak Q_{\mathfrak q}$ of $M$ above it. At that prime choose
a unit squareclass defining the unramified quadratic extension of
the local field $M_{\mathfrak Q_{\mathfrak q}}$; at all the other
primes of $M$ above $\mathfrak q$, choose the square class.
These instructions include every dyadic prime. Such a unit class
exists also in residue characteristic two: lift an irreducible
quadratic $T^2-T+\overline c$ from the residue field and use its
unit discriminant $1-4c$. Its roots generate the unramified quadratic
extension. The prescribed local squareclasses are open, since
sufficiently deep principal units are squares by Hensel's lemma.
Weak approximation therefore gives $b\in M^\times$ that is a unit
at all these primes and has all the prescribed squareclasses.

Choose a $V$-invariant integral ray modulus $\mathfrak m$, supported
on the primes above $S$, deep enough that multiplication by an
element congruent to $1$ modulo $\mathfrak m$ changes none of these
local squareclasses. Let $E/M$ be its ray class field. Invariance
of the modulus makes $E/F$ Galois: every $F$-automorphism of $M$
preserves the defining ray subgroup and hence its class field.
Let $g\in\operatorname{Gal}(E/M)$ be the Artin image of the ray
class of the principal fractional ideal $(b)$.

Chebotarev applied to $E/F$ supplies a prime
$\mathfrak p\notin S$, unramified in $E$, and a prime of $E$ above
it whose Frobenius is $g$; first choose its conjugacy class and then
choose the prime to obtain this representative. Since $g$ acts
trivially on $M$, the prime $\mathfrak p$ splits completely in $M/F$.
For the resulting prime $\mathfrak P$ of $M$, its ray class is
$[(b)]$. By the definition of ray equivalence there is therefore
an element $\eta\in\mathcal O_M$ such that
\[
             \mathfrak P=(\eta),\qquad \eta/b\equiv1\pmod{\mathfrak m}.
\]
Here the equality of ideals makes $\eta$ integral, and the congruence
is only at the modulus primes. This application uses
\cite[Chapter V, Theorem 3.23, p.~164]{C_MilneCFT2020}; it requires
no simultaneous inverse-Galois or embedding theorem.

Now put
\[
       \widetilde N=M\bigl(\sqrt{\sigma\eta}:\sigma\in V\bigr).
\]
The principal ideals $(\sigma\eta)=\sigma\mathfrak P$ are distinct.
Their valuations show that the classes of all $\sigma\eta$ are
linearly independent in $M^\times/M^{\times2}$. Consequently
$\operatorname{Gal}(\widetilde N/M)=\mathbf F_2[V]$, with $e_\sigma$
changing the sign of $\sqrt{\sigma\eta}$ alone. The natural action
of $V$ lifts by the formulas
\[
 g\bigl(\sqrt{\sigma\eta}\bigr)=\sqrt{g\sigma\eta}.
\]
These formulas define an action, not merely separate lifts: the
algebra with these square-root generators is a field by the
independence just proved, and permutation of its labeled generators
respects every defining relation and the multiplication law of $V$.
Thus $\operatorname{Gal}(\widetilde N/F)=\mathbf F_2[V]\rtimes V$.

Use Lemma~\ref{A:A2-wreath-calculation} to quotient by $\ker\phi$.
The resulting $N/F$ has group $G=V\wr C_2$, and the old extension
$M/F$ is the quotient $\Sigma:G\to V$. Let $L=N^H$, where
$H=V\oplus V$; then $[L:F]=2$.

We first check ramification at every prime. At a nondyadic prime
of $M$ outside the primes over $\mathfrak p$, all radicands are
units, so the local quadratic extensions are split or unramified.
At dyadic primes they are split or unramified by the imposed local
squareclasses, which are permuted under $V$. At the prime
$\sigma\mathfrak P$ the only radicand of odd valuation is
$\sigma\eta$. Thus the inertia group there in
$\widetilde N/M$ is exactly $\langle e_\sigma\rangle$; the other
unit radicands contribute only unramified extensions. Because
$\mathfrak p$ splits in $M/F$, this is also the inertia in
$\widetilde N/F$. In $G$ its image is the complement conjugate
$\langle\phi(e_\sigma)\rangle$, whose intersection with $H$ is
trivial. It follows that $N/L$ is unramified at every finite prime.
There are no real places.

We next verify inertness at each old prime and splitting above two.
Fix $\mathfrak q\in S$ and let $v\in V$ be its Frobenius in $M/F$.
If $v=0$, all the primes above $\mathfrak q$ have degree one. At
one such prime exactly one conjugate radicand has the nonsquare
unramified local class, so its Frobenius in the regular module is
$e_g$ for some $g$. Its image $(1,g)\in A$ has order two and
nonzero image in $G/H$. Thus $L/F$ is unramified inert at
$\mathfrak q$, and $N/L$ is split at the prime over it.
In particular this holds at every dyadic prime, since $M/F$ is
$2$-split.

If $v\ne0$, the decomposition subgroup in $M/F$ is
$\langle v\rangle$, of order two. At a chosen prime above
$\mathfrak q$, precisely two conjugate radicands have the
nonsquare unramified local class, with indices $g$ and $g+v$.
Therefore the Frobenius of $\widetilde N/M$ is $e_g+e_{g+v}$,
and its image under $\phi$ is $(0,v)\in A$.
If $z\in G$ is the Frobenius over $\mathfrak q$, then
$\Sigma(z)=v$ and $z^2=(0,v)$ in $A$. The square formula in
Lemma~\ref{A:A2-wreath-calculation} says that a lift $(b,u;v)$
has square $(0,bv)$; since $v\ne0$, we have $b=1$.
Hence $z$ again has nonzero image in $G/H$, proving unramified
inertness in $L/F$.

Finally, Artin reciprocity gives a surjection
$C_L\twoheadrightarrow H$ because $N/L$ is an unramified,
$2$-split, elementary abelian two-extension. Extension of ideals
corresponds to transfer; the exact compatibility is
\cite[Chapter V, Proposition 3.18 and the paragraph following it,
p.~162]{C_MilneCFT2020}. The transfer calculation gives the
commutative diagram
\[
 \begin{array}{ccc}
 C_F&\longrightarrow&C_L\\
 \big\downarrow\wr&&\big\downarrow\\
 V&\xrightarrow{\ v\mapsto(v,v)\ }&V\oplus V.
 \end{array}
\]
One may compute this diagram on ideals prime to the finite
ramification set of $N/F$; every ideal class has such a representative.
The bottom map is injective, so the top map is injective. The
surjection on the right gives $\dim C_L\ge2\dim V=2r$.
\end{proof}

\begin{theorem}[Refuted; A.2 even for an excellent Dedekind domain]
\label{A:A2-persistent-arithmetic-counterexample}
There exists an excellent Dedekind domain $R$, of characteristic
zero with $2$ invertible, whose fraction field $K$ is algebraic
over $\mathbf Q$ and satisfies
\[
 \operatorname{cd}_2(K)=1,
 \qquad
 \operatorname{cd}_2(U)=2
 \quad\text{for every nonempty affine open }U\subseteq\operatorname{Spec}R.
\]
In fact, $\operatorname{Pic}(U)/2$ is infinite dimensional over
$\mathbf F_2$ for every nonempty open $U$.
\end{theorem}
\begin{proof}
We start with $F_0=\mathbf Q(\sqrt{-35})$ and
$\mathcal O_{F_0}=\mathbf Z[w]$, where
$w=(1+\sqrt{-35})/2$ satisfies $w^2-w+9=0$.
The discriminant is $-35$. Minkowski's bound
$(2/\pi)\sqrt{35}<4$ shows that the ideal class group is
generated by primes of norm $2$ or $3$; we use
\cite[Section 4, Theorem 4.3, pp.~70--71]{C_MilneANT2020}.
The polynomial of $w$ is irreducible modulo $2$, so $2$ is
inert, its prime ideal is $(2)$, and it has norm $4$.
Modulo $3$ the polynomial has the distinct roots $0$ and $1$.
Thus the primes of norm $3$ are
$P=(3,w)$ and $\overline P=(3,w-1)$. Since $N(w)=9$ and
$w$ is a unit at $\overline P$, we have $(w)=P^2$.
Moreover $P$ is not principal: an element of norm $3$ would
give integers $x,y$ with
\[
 3=N(x+yw)=x^2+xy+9y^2
       =\left(x+\frac y2\right)^2+\frac{35}4y^2,
\]
which is impossible if $y\ne0$ and is impossible for $y=0$
as well. Therefore the class group is $C_2$, generated by
$P$. Inverting $2$ changes nothing because its prime is
principal, so $\dim C_{F_0}=1$.

Enumerate the rational primes, starting with $2$. Having chosen
$F_n$, let $S_n$ consist of all its primes above the first $n+1$
rational primes. Apply Lemma~\ref{A:A2-quadratic-step} to obtain
$F_{n+1}/F_n$, and continue. Set
\[
 K=\bigcup_{n\ge0}F_n,
 \qquad R=\bigcup_{n\ge0}\mathcal O_{F_n}[1/2].
\]
Every fixed rational prime is, from some stage onward, unramified
and inert at every prime above it at every subsequent step.
Consequently it has only finitely many primes above it in $K$,
and all ramification indices along such a prime are bounded.
The ring $R$ is the full integral closure of $\mathbf Z[1/2]$
in $K$. Its localizations at nonzero primes are discrete valuation
rings: they are unions of compatible valuation rings, and bounded
ramification makes the value group discrete. Every nonzero element
has only finitely many prime divisors, by finite splitting.

For completeness, these two facts prove noetherianity. If $I$ is
a nonzero ideal, choose $0\ne a\in I$ and list the finitely many
maximal ideals $P_1,\ldots,P_s$ containing $a$. At each $P_j$
choose $b_j\in I$ attaining the least valuation of a nonzero
element of $I$. Then $J=(a,b_1,\ldots,b_s)$ and $I$ have the
same localizations at all maximal ideals, so $I=J$.
Thus $R$ is Dedekind. This is also the finite-splitting and
discrete-valuation criterion in \cite[Theorem and Corollary,
pp.~40--42]{C_PopescuVraciu1985}. A Dedekind domain with fraction
field of characteristic zero is excellent
\cite[Proposition 15.53.3(4), Tag 07QS]{C_StacksArithmetic}.

We check the generic cohomological dimension with arbitrary
coefficients. Every finite extension $T/K$ has the form $EK$
for a number field $E$ containing some $F_n$. A class in
$\operatorname{Br}(T)[2]$ descends to a Brauer class of order
dividing two over such an $E$, after increasing the finite
stage if necessary. Only finitely many local invariants of
this class are nonzero. At a prime $w$ of $E$ lying above a
prime of $F_n$, the local degrees in $EF_m/E$ have unbounded
two-adic valuation, uniformly over all prolongations of $w$.
To see this precisely, first enlarge $n$ so that the underlying
rational prime is among the prescribed primes. Every local
degree in $F_m/F_n$ at that prime is then $2^{m-n}$. The tower
degree identity in a common completion gives
\[
 v_2\bigl([\,(EF_m)_{w_m}:E_w\,]\bigr)
 \ \ge\ m-n-v_2\bigl([\,E_w:(F_n)_v\,]\bigr).
\]
Thus a single sufficiently large $m$ makes all local restriction
degrees even at the finite support of the Brauer class.
All archimedean places are complex. Restriction multiplies
each local invariant by its local degree, and the
Brauer--Hasse--Noether injection then shows that the class
vanishes over $EF_m$. We have proved
\[
            \operatorname{Br}(T)[2]=0
                 \quad\text{for every finite extension }T/K.
\]
The precise local-global and restriction inputs are
\cite[Chapter III, Proposition 1.8, p.~100; Chapter VII,
Theorem 7.1 and Proposition 7.2, pp.~219--220; Chapter VIII,
Theorem 4.2, p.~241]{C_MilneCFT2020}.

As an algebraic extension of a totally imaginary number field,
$K$ initially has $\operatorname{cd}_2(K)\le2$; we use
\cite[Proposition 8.3.18, with all primes in $S$]{A_NSW}
and closed-subgroup monotonicity. For a finite discrete
$\mathbf F_2[\operatorname{Gal}(K^{\mathrm{sep}}/K)]$-module $M$,
choose a finite Galois extension $T/K$ trivializing it.
A finite sum of regular permutation modules surjects onto $M$.
Its second cohomology vanishes by Shapiro and
$H^2(T,\mathbf F_2)=H^2(T,\mu_2)=\operatorname{Br}(T)[2]=0$.
The kernel has zero third cohomology by the preliminary bound two.
The long exact sequence gives $H^2(K,M)=0$. Filtered colimits
and torsion-coefficient d\'evissage give $\operatorname{cd}_2K\le1$.
There are nonzero primes of $R$, by lying over any nondyadic
prime of $\mathbf Z[1/2]$. A uniformizer in the discrete
valuation ring $R_P$ has odd valuation in $K$, so it is not a
square. Kummer theory gives a nonzero class in $H^1(K,\mu_2)$.
Thus $\operatorname{cd}_2K=1$.

Invertible modules and their isomorphisms descend through
filtered ring colimits, and taking the cokernel of multiplication
by two commutes with such colimits. Hence
\[
 \operatorname{Pic}(R)/2
     =\varinjlim_n \operatorname{Pic}(\mathcal O_{F_n}[1/2])/2.
\]
All transition maps are injective, and the dimensions are at least
$2^n$, by Lemma~\ref{A:A2-quadratic-step}. The resulting vector
space is therefore infinite dimensional. Every nonempty open
of the noetherian integral curve $\operatorname{Spec}R$ removes
only finitely many closed points. Its Picard group is the quotient
of $\operatorname{Pic}(R)$ by the classes of those points.
After reduction modulo two this kills at most a finite-dimensional
subspace. Consequently $\operatorname{Pic}(U)/2$ remains infinite
dimensional for every such $U$.

The Kummer sequence injects $\operatorname{Pic}(U)/2$ into
$H^2(U,\mu_2)$, so every nonempty affine open has
$\operatorname{cd}_2(U)\ge2$. The generic field is nonorderable,
and every closed-point residue field has positive characteristic.
Corollary~\ref{A:generic-plus-dimension},
the generic-field-plus-dimension bound for noetherian schemes,
gives the reverse
inequality $\operatorname{cd}_2(U)\le\operatorname{cd}_2K+1=2$.
This proves both the asserted exact dimensions and the failure
of the full eventual-equality quantifiers of A.2.
\end{proof}

\paragraph{Finite verification and scope.}
The accompanying exact-arithmetic program
\nolinkurl{parts/C/code/check_quadratic_tower.py} checks the wreath-group
identities for $\dim V=1,2,3$, the transfer from its coset definition,
every local Frobenius-square implication, the inertia complements,
and the initial norm calculation and reduced forms. Its saved output
\nolinkurl{parts/C/output/quadratic_tower_certificate.json} records
$111{,}784$ passed assertions in $0.068522113$ seconds. These finite checks
supplement the proofs above; the infinite tower is justified by the
quadratic-step lemma, Chebotarev, and the explicit noetherianity and
Brauer arguments. No positive Iwasawa $\mu$-invariant or finite-stage
class-rank growth without persistence is assumed.

\begin{corollary}[Refuted; persistent one-dimensional forms of A.1, A.6, L1 and L2]
\label{A:A2-sharp-curve-consequences}
For the excellent Dedekind domain in
Theorem~\ref{A:A2-persistent-arithmetic-counterexample}, every nonempty
affine open $U$ satisfies
\[
 \sup_{x\in U}\bigl(\cd_2\kappa(x)+\dim\mathcal O_{U,x}\bigr)=1,
 \qquad \cd_2(U)=2.
\]
Every residue field is nonorderable. Thus the failures of A.1 and
both statements of A.6 persist on every nonempty affine open, with
the smallest positive generic cohomological dimension. The upper
curve bound $\cd_p(R)\leq\cd_p(K)+1$ cannot in general be improved
by any nonempty localization, even for excellent regular domains.
\end{corollary}
\begin{proof}
A closed-point residue field is the union of finite fields whose
successive degrees are eventually two, because each retained prime
is eventually inert in every quadratic step. Its absolute Galois
group is a closed subgroup of $\widehat{\mathbf Z}$ with trivial
pro-$2$ subgroup, so its $2$-cohomological dimension is zero.
Its local ring has dimension one. At the generic point the two
terms are respectively one and zero. These exhaust the points of
the Dedekind scheme, proving the displayed supremum. The preceding
theorem gives $\cd_2(U)=2$, and gives the same value for every
nonempty affine open inside $U$; this also computes the supremum
asked for in A.6(ii). The generic field contains the totally
imaginary number field $\mathbf Q(\sqrt{-35})$, and the closed
residue fields have positive characteristic, so none is orderable.

The proved replacement retains the generic lower witness of
Proposition~\ref{A:spread} and the sharp curve upper bound of
Corollary~\ref{A:generic-plus-dimension}. Theorems
\ref{A:Dedekind-Milnor-criterion} and \ref{A:curve-normalization}
give the exact additional condition for eventual equality on
every noetherian integral curve, without excellence or normality:
one localization must make the simultaneous Milnor residue maps
surjective on every finite separable normalization. The constructed
domain fails that condition already for its own Picard group.
\end{proof}


\subsection{A.3 and A.4: a mixed-characteristic counterexample}

\begin{theorem}[Refuted; A.3 and the upper bound in A.4]
\label{A:mixedreal}
There is an excellent semilocal integral ring $A$ for which $3$ is a unit, $2$ is nonzero and not a unit,
\[
\cd_2(A)=+\infty,\qquad
\cd_2(\Frac A)=3,\qquad \delta(V(2))=1.
\]
Thus both the unnumbered conjecture on 31/085 and the bound
\[
\cd_2(A)\leq\max\bigl(\cd_2(\Frac A),\delta(V(2))+3\bigr)
\]
fail already for quasi-excellent schemes.
\end{theorem}
\begin{proof}
Put $T=\Z[x,y]/(x^2+y^2)$. This is a domain: its defining polynomial is primitive and irreducible over $\Q$. Consider the incomparable prime ideals
\[
P=(x,y),\qquad Q=(2,x-1,y-1),
\]
and the multiplicative set $S=T\setminus(P\cup Q)$. Set $A=S^{-1}T$. Prime avoidance shows that the maximal ideals of $A$ are precisely $PA$ and $QA$. The residue field at $PA$ is $\Q$ and that at $QA$ is $\F_2$; equivalently, these are the fraction fields of $T/P$ and $T/Q$ after the semilocalization. The ring $A$ is excellent, since $T$ is of finite type over $\Z$ and excellence is stable under localization. The integer $3$ belongs to $S$, whereas $2\in Q$ and is nonzero in the domain $A$.

The fraction field is $\Q(i)(x)$, by the identity $i=y/x$ and $i^2=-1$. The number field $\Q(i)$ is totally imaginary and has $2$-cohomological dimension two by \citet[Theorem I.4.10(c)]{A_Milne2006}. The transcendence-degree formula \citet[Exposé X, Théorème 2.1]{A_SGA4} therefore gives $\cd_2(\Frac A)=3$.

The reduction of $T/2T$ is $\F_2[x,y]/(x+y)\simeq\F_2[x]$. Its unique generic point survives the localization: the corresponding prime of $T$ is contained in $Q$, and hence is disjoint from $S$. The sole maximal point of $V(2)$ thus has residue field $\F_2(x)$, whose $2$-rank is one. Nilpotents in $T/2T$ do not alter its points or residue fields. Consequently $\delta(V(2))=1$.

Finally, $\Spec\Q\hookrightarrow\Spec A$ is closed. The orderable field $\Q$ has infinite $2$-cohomological dimension. For example, its absolute Galois group contains the order-two subgroup attached to a real closure, which already has infinite $2$-cohomological dimension. Lemma~\ref{A:basic} gives $\cd_2(A)=+\infty$. The proposed finite upper bound would be $\max(3,4)=4$.
\end{proof}

\subsection{A.8 and A.9: the local lower bound without excellence}

\begin{lemma}[Proved from known valuation formulas]
\label{A:valuation-lower}
Let $F$ be a discretely valued field, $V$ its valuation ring, and $E$ its residue field. If $\ell=2$, suppose that $E$ is not orderable. Then
\[
d_\ell(F)\geq1+d_\ell(E).
\]
Neither $V$ nor $F$ is assumed excellent or henselian.
\end{lemma}
\begin{proof}
Let $\widehat F$ be the completion. When $\operatorname{char}F\neq\ell$, Krasner approximation identifies the absolute Galois group of the henselization with that of the completion, and embeds this group as a closed subgroup of $G_F$. Therefore $\cd_\ell(F)\geq\cd_\ell(\widehat F)$; this is the argument of \citet[\S3.2]{A_GO2008}, using \citet[Exposé X, Lemme 2.2.1]{A_SGA4}. If $\operatorname{char}F=\ell$, the corresponding inequality $d_\ell(F)\geq d_\ell(\widehat F)$ is \citet[\S3.3, the displayed lemma]{A_GO2008}. That lemma expressly assumes only a discrete valuation.

The complete discrete valuation ring $\widehat V$ is excellent and henselian. If $\operatorname{char}E=\ell$, Kato's theorem in \citet[\S2.3, Théorème 2.3.1]{A_GO2008} gives
$d_\ell(\widehat F)=1+d_\ell(E)$. If $\operatorname{char}E\neq\ell$, the same equality for ordinary cohomological dimension is \citet[Exposé X, Théorème 2.2(i)]{A_SGA4}. The nonorderability assumption removes the familiar finite-extension exception at $\ell=2$. Combining these comparisons proves the assertion, including infinite values.
\end{proof}

\subsection{A regular excellent affine counterexample without real points}

\begin{theorem}[Known construction; consequences for A.1, A.5, A.6, A.7]
\label{A:nagata}
Fix a prime $\ell$. There exist a strictly henselian regular noetherian local $\overline{\Q}$-algebra $B$ of dimension two and a nonzero element $F$ of its maximal ideal such that, for $U=\Spec B[1/F]$ and $K=\Frac B$,
\[
\dim U=1,\qquad \cd_\ell(U)=3,\qquad \cd_\ell(K)=2.
\]
The ring $B$ is not quasi-excellent, but the affine scheme $U$ is excellent and regular. Every residue field in this construction is nonorderable.
\end{theorem}
\begin{proof}
We specify the Nagata--Gabber construction of \citet[Exposé XVIII$_B$, \S2, pp.~464--467]{A_ILO2014}. Put $k=\overline{\Q}$, $W=k[t]_{(t)}$, $K_0=k(t)$, and
\[
\varphi=\exp(t)-1=\sum_{n\geq1}\frac{t^n}{n!}\in t\,k[[t]].
\]
This is transcendental over $K_0$. Indeed an algebraic relation would give a nontrivial finite relation $\sum_{j=0}^m P_j(t)e^{jt}=0$ with $P_j\in\C[t]$. For the largest $j$ with $P_j\neq0$, dividing by $e^{jt}$ and letting the real variable $t$ tend to $+\infty$ contradicts polynomial growth. An identity of formal series here is an identity of entire functions, so this analytic argument proves the formal assertion.

Let $f=(y-\varphi)^\ell\in\widehat W[y]$. Choose $f_n\in W[y]$ with $f_n\equiv f\pmod{t^n}$, and let $F$ now denote an indeterminate. Define the explicitly specified subring
\[
\mathcal D=
\bigcup_{n\geq1}W[y]\bigl[(F-f_n)/t^n\bigr]
\ \subset K_0[y,F].
\]
The inclusions between successive terms follow from $f_{n+1}-f_n\in t^nW[y]$. Evaluating $F$ at $f$ gives an injection $\mathcal D\hookrightarrow\widehat W[y]$; injectivity follows from the transcendence of $\varphi$. Let $A$ be the localization of $\mathcal D$ at the inverse image of $(t,y)\subset\widehat W[y]$, and put $B=A^h$.

For completeness, the algebraic verifications supplied in the opened source have the following precise roles. The map $\mathcal D\to\widehat W[y]$ is flat: modulo $t$ it is an isomorphism, while after inverting $t$ it is the flat map
\[
K_0[y,F]\longrightarrow\widehat K_0[y],\qquad
F\longmapsto(y-\varphi)^\ell.
\]
This factors through $K_0[y,F']$, $F\mapsto(F')^\ell$, followed by $F'\mapsto y-\varphi$. For the latter map, change coordinates to $Z=y-F'$: it becomes $K_0[y,Z]\to\widehat K_0[y]$, with $Z\mapsto\varphi$. It is flat because it factors through the field embedding $K_0(Z)\hookrightarrow\widehat K_0$ and then adjoins $y$. After the indicated localization the flat map is local and faithfully flat. The source is therefore noetherian, and its completion is $k[[t,y]]$. Hence $A$, and then $B$, are regular of dimension two. Their residue field is $k$, so the henselian ring $B$ is strictly henselian. The source's \S2.2 proves that $A/F$ is integral and regular away from the closed point. Its \S2.3 proves that $B/F$ is still integral, by comparison of the connected components of the punctured divisor with its completion. In $\widehat B=k[[t,y]]$ the element $F$ becomes
\[
F=(y-\varphi)^\ell.
\]
Thus the completion of the reduced ring $B/F$ is nonreduced, and $B$ cannot be quasi-excellent.

We now give the cohomological calculation of \S2.4 with the maps needed here. Let $X=\Spec B$, let $m$ be its closed point, put $T=X\setminus\{m\}$, and let $C=V(F)$. The regular one-dimensional scheme $C\setminus\{m\}$ is the spectrum of the fraction field of the one-dimensional strictly local domain $B/F$. Its $\ell$-dimension is one by Theorem~\ref{A:strict-field}. The localization sequence and absolute purity therefore give an exact sequence
\[
H^1(C\setminus\{m\},\Lambda(-1))
\xrightarrow{\mathrm{Gys}_F}H^3(T,\Lambda)
\longrightarrow H^3(U,\Lambda)\longrightarrow0,
\quad \Lambda=\Z/\ell.
\]
The comparison maps from the punctured schemes to their completions are isomorphisms in this sequence; the precise input is the completion comparison cited and applied in \S2.4 of Exposé XVIII$_B$. On the completion the Gysin map for $F=(y-\varphi)^\ell$ is $\ell$ times the Gysin map for $y-\varphi$. It is therefore zero with $\Lambda$ coefficients. Compatibility of Gysin maps with the comparison isomorphisms shows that $\mathrm{Gys}_F$ itself is zero. Purity at the closed point of the regular strictly local surface gives
\[
H^3(T,\Lambda)=\Lambda(-2),
\qquad H^3(U,\Lambda)=\Lambda(-2)\neq0.
\]
The general upper bound for opens of a two-dimensional strictly local noetherian scheme is three, by \citet[Exposé XVIII$_A$, Théorème 1.1]{A_ILO2014}. Hence $\cd_\ell(U)=3$. Theorem~\ref{A:strict-field} gives $\cd_\ell(K)=2$.

Finally, $U$ is a nonempty proper open of the regular local surface. It has no point of height two, while the height-one prime $(t)$ belongs to it because $F\bmod t=y^\ell\neq0$. Thus $\dim U=1$. Its coordinate ring is a noetherian regular domain of dimension one, hence a Dedekind domain, and its fraction field has characteristic zero. Such a Dedekind domain is excellent by \citet[Tag 07QS, Proposition 15.53.3(4)]{A_Stacks}. This last step concerns $U$ itself, and must not be confused with the nonexcellence of the larger local ring $B$. Since $k\subset B$ contains a square root of $-1$, all residue fields are nonorderable.
\end{proof}

\begin{corollary}[Refuted in both requested classes: A.1 and A.6(i)--(ii)]
\label{A:nagata-global-consequences}
A.1 fails even for excellent regular affine schemes without real residue fields. Both assertions of 31/082, A.6(i) and A.6(ii), also fail for such schemes. More precisely, for the excellent affine scheme $U$ of Theorem~\ref{A:nagata},
\[
\sup_{u\in U}\bigl(\cd_\ell\kappa(u)+\dim\OO_{U,u}\bigr)=2
\quad\text{but}\quad \cd_\ell(U)=3.
\]
\end{corollary}
\begin{proof}
At the generic point the summand is $2+0=2$. Every other point of $U$ is given by a height-one prime $\mathfrak p$ of $B$, different from $(F)$. Its residue field is $\Frac(B/\mathfrak p)$. The quotient $B/\mathfrak p$ is a one-dimensional strictly henselian local domain, so Theorem~\ref{A:strict-field} makes its fraction field's $\ell$-dimension one. The local ring $\OO_{U,u}=B_{\mathfrak p}$ also has dimension one. Thus every summand is two. The bound in A.6(i) would say $3\leq2$. The supremum in A.6(ii) includes the affine open $U$ itself and is consequently at least three, so cannot be two. The same scheme refutes A.1 by $3>2$.
\end{proof}

\begin{corollary}[Refuted without quasi-excellence: A.5 and affine A.7]
\label{A:nagata-local-consequences}
Without quasi-excellence, the upper bound in A.5 is false even in residue characteristic zero. A.7(i) and the affine clause of A.7(iii) are likewise false without quasi-excellence of the strictly local base.
\end{corollary}
\begin{proof}
In A.5 take the ring $B$ of Theorem~\ref{A:nagata} and its affine open $U$. We have $\dim B=2$, $\delta(k)=0$, and $\cd_\ell(k)=0$, whereas $\cd_\ell(U)=3$.

For A.7 use $Y=\Spec B$ and $X=U$. The open immersion is affine and of finite type, the closed fibre is empty, and $X'=X$ has dimension one. Thus the claimed bound in (i) is two. The morphism is dominant with zero-dimensional generic fibre, so the affine bound in (iii) is also two. Both contradict the value three. The excellence of the \emph{source} $U$ does not satisfy the missing excellence hypothesis on the \emph{base} $Y$.
\end{proof}

\subsection{A.5: proved positive cases and Artin's N.B.}

\begin{proposition}[Known, with descent made explicit]
\label{A:hens-affine-unit}
Let $(A,\mathfrak m,k)$ be a henselian quasi-excellent noetherian local ring of dimension $d$, let $\ell$ be invertible in $A$, and let $U\subset\Spec A$ be affine. Then
\[
\cd_\ell(U)\leq d+\cd_\ell(k).
\]
Without quasi-excellence, for $d>0$ the proved replacement is
$\cd_\ell(U)\leq2d-1+\cd_\ell(k)$. The example of Theorem~\ref{A:nagata} attains the latter value when $d=2$ and $k$ is algebraically closed.
\end{proposition}
\begin{proof}
Let $A^{\mathrm{sh}}$ be a strict henselization and $\widetilde U=U\times_A A^{\mathrm{sh}}$. The pro-étale Galois covering $\Spec A^{\mathrm{sh}}\to\Spec A$, with group $G_k$, gives the Hochschild--Serre spectral sequence
\[
H^a\bigl(G_k,H^b(\widetilde U,F|_{\widetilde U})\bigr)
\ \Longrightarrow\ H^{a+b}(U,F)
\]
for every $\ell$-torsion sheaf $F$; this is \citet[Exposé VIII, \S8]{A_SGA4}. Strict henselization preserves dimension and quasi-excellence. The affine theorem of \citet[Exposé XV, Corollaire 1.2.2]{A_ILO2014} gives vanishing for $b>d$. Without quasi-excellence, Exposé XVIII$_A$, Théorème 1.1 gives vanishing for $b>2d-1$. In either case Galois cohomology vanishes for $a>\cd_\ell(k)$. The indicated total-degree bounds follow. The case $d=0$ is simply the field case, allowing nilpotents, which do not change the étale topos.
\end{proof}

\begin{corollary}[Known; A.5.N.B. on 31/087]
\label{A:artin-NB}
For a henselian quasi-excellent local domain with residue characteristic zero, the bound in A.5 holds for every prime, with $\delta(k)=0$. Artin's original result already proves this case.
\end{corollary}
\begin{proof}
Every prime is a unit in such a local ring, so Proposition~\ref{A:hens-affine-unit} applies. To identify the historical input exactly, a henselian quasi-excellent local ring is excellent: \citet[Exposé I, Corollaire 6.3(ii)]{A_ILO2014} proves that a henselian local ring with geometrically regular formal fibres is excellent. Equivalently, its integral quotients have integral completions, and the formal-catenarity criterion applies. The ring has equal characteristic zero. Therefore \citet[Exposé XIX, Corollaire 6.2, deduced from Théorème 6.1 bis]{A_SGA4} gives directly
$\cd_\ell(U)\leq\dim A+\cd_\ell(k)$ for every prime $\ell$. Artin's resolution input in that exposé is available in characteristic zero. This justifies the N.B. in precisely the quasi-excellent case, while Corollary~\ref{A:nagata-local-consequences} shows why the unrestricted noetherian version fails.
\end{proof}

\begin{proposition}[Known; the generic-characteristic open in A.5]
\label{A:wild-generic-open}
Let $A$ be a henselian quasi-excellent local domain of mixed characteristic $(0,p)$, with residue field $k$. For every nonempty affine open $U\subset\Spec A[1/p]$,
\[
\cd_p(U)=\dim A+d_p(k)
\leq\dim A+\delta(k)+\cd_p(k).
\]
\end{proposition}
\begin{proof}
As in Corollary~\ref{A:artin-NB}, the henselian quasi-excellent ring is excellent. The equality is precisely \citet[Théorème 6.1]{A_GO2008}, whose hypothesis places the entire affine open in $\Spec A[1/p]$.

For the inequality, $d_p(k)$ is either $\delta(k)$ or $\delta(k)+1$ when the $p$-rank is finite; this follows at once from its definition by differential forms and logarithmic $H^1$. If $\cd_p(k)=1$, this proves the desired estimate. If $\cd_p(k)=0$, every $p$-primary discrete Galois module has zero positive cohomology, also after finite extensions. The logarithmic differential sheaves used in the definition of $d_p$ are such modules, so $d_p(k)=\delta(k)$. Infinite $p$-rank makes the displayed upper bound vacuous. We do not apply this theorem to an open meeting the characteristic-$p$ fibre.
\end{proof}

\subsection{A.5: a wild counterexample on a regular complete local surface}

\begin{theorem}[Refuted; A.5 even for excellent regular complete local rings]
\label{A:wild-affine-counterexample}
Let $k=\overline{\F}_2$, let $W=W(k)$ be its ring of Witt vectors, and put
\[
 A=W[[t]],\qquad U=\Spec A[1/t].
\]
Then $A$ is an excellent regular henselian local domain of dimension two,
its fraction field has characteristic zero, and
\[
 \cd_2(U)=3>2=\dim A+\delta(k)+\cd_2(k).
\]
In fact $H^3(U,\F_2)\neq0$. Thus the residual-characteristic case of
Conjecture 4 on 31/085--087 is false even with quasi-excellence.
\end{theorem}
\begin{proof}
The ring $W$ is a complete discrete valuation ring with uniformizer $2$
and residue field $k$. Hence $A$ is a complete regular local domain with
regular parameters $(2,t)$. Complete noetherian local rings are excellent.
The residue field is algebraically closed and perfect, so both terms
$\delta(k)$ and $\cd_2(k)$ are zero.

Write
\[
 R=\Spec A[1/2],\quad V=\Spec A[1/(2t)],\quad
 B=k[[t]],\quad E=k((t)),\quad C=\Frac W.
\]
Let $D=(A_{(2)})^h$ be the henselization of the discrete valuation ring
$A_{(2)}$, and let $L=\Frac D$. Its residue field is $E$.
All symbols below are degree-two Kummer symbols with coefficients
$\mu_2^{\otimes2}=\F_2$ on characteristic-zero schemes.

The complement $z=U-V$ is the single closed point $\Spec E$.
Henselian invariance gives $\cd_2(D)=\cd_2(E)\leq1$.
Gabber--Orgogozo gives $d_2(E)=1$ and $\cd_2(L)=2$:
apply their Th\'eor\`eme 1.2 first to $k[[t]]$ and then to the excellent
henselian discrete valuation ring $D$.
By Proposition~\ref{A:wild-generic-open}, both $R$ and $V$ have
$2$-cohomological dimension two.
\'Etale excision for henselization, followed by the localization sequence
for the closed point of $\Spec D$, gives
\[
 H^3_z(U,\F_2)=H^3_z(\Spec D,\F_2)
       \simeq H^2(L,\F_2).
\]
Consequently localization on $U$ gives an exact sequence
\begin{equation}\label{A:wild-cokernel}
 H^2(V,\F_2)\longrightarrow H^2(L,\F_2)
       \longrightarrow H^3(U,\F_2)\longrightarrow0.
\end{equation}
For arbitrary $2$-primary torsion coefficients, the same localization
sequences give vanishing above degree three: the terms on $D$, $L$ and
$V$ vanish above degrees one, two and two respectively. Thus
$\cd_2(U)\leq3$.

We now exhibit an element outside the image in
\eqref{A:wild-cokernel}. The only point of $R-V$ is $(t)$, with residue
field $C$. Purity applies here because $2$ is invertible on $R$.
The localization sequence and $H^3(R,\F_2)=0$ show that
\[
 H^2(R,\F_2)\longrightarrow H^2(V,\F_2)
       \longrightarrow H^1(C,\F_2)\longrightarrow0
\]
is exact. Every element of the final group is represented by some
$u\in C^\times$ and lifted by the symbol $\{t,u\}$, since its residue
at $(t)$ is the Kummer class of $u$. Therefore every class coming from
$H^2(V,\F_2)$ has the form
\begin{equation}\label{A:wild-image-form}
 a+\{t,u\},\qquad a\in H^2(R,\F_2),\quad u\in C^\times.
\end{equation}

Here are the precise ramification-filtration inputs. Let
$U^mH^2(L,\F_2)$ be generated by symbols whose first entry lies in
$1+2^mD$ for $m\geq1$, and put $U^0=H^2$.
Since the absolute ramification index is one, the upper critical index
is $2$. The formulas of
\citet[Proposition 6.1.2(i),(iii)]{A_Suzuki2024}
(Kato's symbol filtration) identify
\[
 \operatorname{gr}^0H^2(L,\F_2)=E^\times/E^{\times2},\qquad
 \operatorname{gr}^1H^2(L,\F_2)=\Omega^1_E,
\]
with $v\mapsto\{\widetilde v,2\}$ in degree zero and
$x\,d\log y\mapsto\{1+2\widetilde x,\widetilde y\}$ in degree one.
The restriction $H^2(R,\F_2)\hookrightarrow H^2(L,\F_2)$ is injective
and strict for this filtration; its degree-zero and degree-one images
are respectively
\[
 B^\times/B^{\times2}\subset E^\times/E^{\times2},\qquad
 \Omega^1_B\subset\Omega^1_E.
\]
These are exactly
\citet[Propositions 8.1.4, 8.2.1(i),(iii), and 8.2.3]{A_Suzuki2024}.
All their hypotheses hold: $A$ is regular, excellent and henselian,
$\zeta_2=-1\in A$, the unique prime dividing $2$ is $(2)$, and
$A/(2)=B$ is regular.

Suppose a class of the form \eqref{A:wild-image-form} belongs to
$U^1H^2(L,\F_2)$. We may write $u=2^e u_0$ with $u_0\in W^\times$.
Since $k$ is perfect, dividing $u_0$ by a square of a Teichm\"uller
unit puts it in $1+2W$ without changing its Kummer class.
The degree-zero image of $\{t,u\}$ is $t^e$ modulo squares.
The degree-zero image of $a$ is a unit of $B$ modulo squares.
Its $t$-adic valuation is even. Cancellation in degree zero therefore
forces $e$ to be even. Removing this square factor makes
$\{t,u\}\in U^1$, and strictness forces $a\in U^1$ as well.
Write $u=1+2b$ with $b\in W$. Its degree-one image is
$\bar b\,d\log t$ (the sign is immaterial in characteristic two).
Thus the degree-one image of every class in
\[
 \operatorname{im}\bigl(H^2(V,\F_2)\to H^2(L,\F_2)\bigr)\cap U^1
\]
belongs to
\begin{equation}\label{A:wild-form-obstruction}
 \Omega^1_B+k\,d\log t=B\,dt+k\,t^{-1}dt.
\end{equation}
On the other hand,
\[
 c=\{1+2/t,t\}\in U^1H^2(L,\F_2)
\]
has degree-one image $t^{-1}d\log t=t^{-2}dt$, which is not in
\eqref{A:wild-form-obstruction}. To see that no differential-form
qualification is hidden here, $\{t\}$ is a $2$-basis of both $B$ and
$E$, as $k$ is perfect. Hence their absolute differential modules are
$B\,dt$ and $E\,dt$. The coefficient of $t^{-2}$ is an explicit
linear functional separating the displayed form from the displayed
subspace. Thus $c$ is not in the image in \eqref{A:wild-cokernel}, and
its boundary is a nonzero class in $H^3(U,\F_2)$.
\end{proof}

\begin{remark}[Source and scope]
Gabber's 2015 IH\'ES lecture, at 33:17--33:24, pointed us toward a
possible extra degree for affine opens meeting the characteristic-$p$
locus. We used its automatic transcript only as a search lead.
The proof above uses the opened written symbol calculations and exact
localization sequences, not that transcript as a theorem.
No priority is claimed for this example or its derivation.
\end{remark}

\begin{theorem}[Proved under stated hypotheses; a uniform mixed-characteristic replacement]
\label{A:wild-affine-upper}
Let $A$ be an excellent henselian local domain of mixed characteristic
$(0,p)$, of dimension $d$, with residue field $k$. For every affine open
$U\subset\Spec A$,
\[
 \cd_p(U)\leq2d-1+d_p(k)
       \leq2d-1+\delta(k)+\cd_p(k).
\]
For $d=1$ this is the fraction-field bound. For $d=2$ the first bound
is attained by Theorem~\ref{A:wild-affine-counterexample}. If
$U\subset D(p)$ the sharper equality in
Proposition~\ref{A:wild-generic-open} still applies.
\end{theorem}
\begin{proof}
Write $s=d_p(k)$ and $r=\delta(k)$. Infinite $r$ makes the assertion
vacuous; assume $r<\infty$, so $r\leq s\leq r+1$.
The inequality $s\leq r+\cd_p(k)$ was proved in
Proposition~\ref{A:wild-generic-open}.
If $U$ contains the closed point it is $\Spec A$, whose cohomology is
that of $k$ by henselian invariance. We therefore assume that $U$
misses the closed point.

We induct on $d$, proving vanishing for all $p$-primary torsion sheaves.
For $d=1$, a nonempty such $U$ is the fraction field, and
Gabber--Orgogozo Th\'eor\`eme 1.2 gives $\cd_p(U)=1+s$.
For $d=2$, the closed subset $Z=U\cap V(p)$ is a finite set of closed
points. For $z\in Z$, its residue field is the fraction field of the
one-dimensional excellent henselian local domain $A/\mathfrak p_z$.
The equal-characteristic part of the same theorem gives
$d_p(\kappa(z))=1+s$.
Normalize the henselization of $\mathcal O_{U,z}$. It is a finite
product of henselian discrete valuation rings, with residue fields
finite over $\kappa(z)$. The generic fields therefore have
$p$-cohomological dimension $2+s$, by the discrete-valuation case and
finite-extension invariance of $d_p$.
The henselian local ring itself has $p$-cohomological dimension at most
one. Localization and henselian excision give
$H^q_z(U,F)=0$ for $q>3+s$. On $U[1/p]$ the bound is $2+s$ by
Gabber--Orgogozo Th\'eor\`eme 6.1. The localization sequence proves
$\cd_p(U)\leq3+s$.

Assume now $d\geq3$ and put $N=2d-1+s$.
We spell out the support argument, adapting the proof of
\citet[Expos\'e X, Th\'eor\`eme 4.1]{A_SGA4}.
By induction, an affine open in an integral closed subscheme
$\Spec A/\mathfrak q$ of characteristic zero and dimension $e<d$
has $p$-cohomological dimension at most $2e-1+s$.
The same estimate holds for a finite union of such closed subschemes
of dimensions at most $e$. Indeed, the exact sequence for a union of
two closed subsets reduces this assertion to their intersection,
whose characteristic-zero components have smaller dimension and
whose characteristic-$p$ part is affine with $p$-cohomological
dimension at most one. For $e=1$ those intersections lie at the closed
point of $\Spec A$, which $U$ misses; for $e\geq2$ the possible extra
degree from the intersection is at most $2\leq2e-1+s$.
Nilpotents do not change the \'etale topos.

It follows that any sheaf zero at the generic point of $U$ has
cohomological dimension at most $N-2$: approximate it by constructible
subsheaves and cover each support by finitely many integral closed
subschemes. Characteristic-zero components have dimension at most
$d-1$; the characteristic-$p$ part has bound one. The closed-union
exact sequence gives at most $\max(N-2,2)=N-2$ since $d\geq3$.

Let $i:\Spec K\to U$ be the generic inclusion and let $G$ be a finite
$p$-torsion Galois module over $K=\Frac A$.
At a characteristic-zero point $y\in U$ of height $h$, the stalks of
$R^b i_*G$ vanish for $b>h$, by
Theorem~\ref{A:strict-field} applied to the components of the strict
localization. At a characteristic-$p$ point $y$ of height $h$, we have
\[
 \delta(\kappa(y))=d-h+r.
\]
For this identity, apply
\citet[Lemmes 2.1.4--2.1.6 and the Cohen-structure calculation in \S4]{A_GO2008}
to the excellent henselian local domain $A/\mathfrak p_y$:
completion preserves the finite $p$-rank and the completed
fraction field has $p$-rank $\dim(A/\mathfrak p_y)+r$.
Excellence gives $\dim(A/\mathfrak p_y)=d-h$.
Strict henselization preserves $p$-rank. The fraction fields of its
components have cohomological dimension at most
$h+\delta(\kappa(y))=d+r$ by Gabber--Orgogozo.
Thus $R^b i_*G=0$ for $b>d+r$.

If $1\leq b\leq d-1$, every constructible subsheaf of $R^b i_*G$
is supported on the union of $U\cap V(p)$ and finitely many
characteristic-zero closed subschemes of dimensions at most $d-b$.
The closed-union sequence, the induction hypothesis and the affine
characteristic-$p$ bound give
\[
 \cd_p\bigl(U,R^b i_*G\bigr)
    \leq\max\{2(d-b)-1+s,\,2\}\leq N-b-1.
\]
Here the notation means the upper vanishing degree for that sheaf;
filtered-colimit compatibility removes constructibility.
If $d\leq b\leq d+r$, its support lies entirely in the closed affine
characteristic-$p$ subset, so its cohomology vanishes above degree one.
This is sufficient because
\[
 N-b-1\geq d-2+s-r\geq1.
\]

In the Leray spectral sequence
$H^a(U,R^b i_*G)\Rightarrow H^{a+b}(K,G)$, all possible incoming
differentials to $H^q(U,i_*G)$ for $q>N$ vanish by these inequalities.
The abutment vanishes too, since $\cd_p(K)=d+s\leq N$.
Hence $H^q(U,i_*G)=0$ for $q>N$.
Finally, for any constructible $p$-torsion sheaf $F$ on $U$, the kernel
and cokernel of $F\to i_*i^*F$ are zero at the generic point. Their
bound $N-2$ just proved and the two associated short exact sequences
give $H^q(U,F)=0$ for $q>N$. Coefficient d\'evissage and filtered
colimits finish the proof for all $p$-primary torsion sheaves.
\end{proof}


\subsection{A.7: nonseparated schemes change the answer}

\begin{theorem}[Refuted already over an excellent trait]
\label{A:double-trait}
As stated without separatedness, A.7(ii) and the nonaffine clause of A.7(iii) are false, even when the strictly local base is excellent and regular of dimension one.
\end{theorem}
\begin{proof}
Let $k$ be algebraically closed of characteristic different from $\ell$, let $A=k[[t]]$, and put $Y=\Spec A$, $\eta=\Spec k((t))$. Glue two copies $Y_1,Y_2$ of $Y$ along their identical open subschemes $\eta=D(t)$. Call the resulting scheme $X$, with its evident morphism $f:X\to Y$.

The morphism is locally of finite type and quasi-compact, hence of finite type: it is covered by the two affine opens $Y_i$, each mapping isomorphically to $Y$. The overlap is affine and quasi-compact. The scheme is reduced and has a common dense generic point, hence is integral. It is nonseparated and therefore not affine. Its closed fibre consists of two reduced $k$-points, so $\dim X_0=0$; its open part $X'$ is $\eta$, so $\dim X'=0$. The generic fibre of $f$ has dimension $d=0$, and $f$ is dominant.

Write $\Lambda=\Z/\ell$. Each $Y_i$ is strictly local, so $H^q(Y_i,\Lambda)=0$ for $q>0$. Étale Mayer--Vietoris for this Zariski open cover gives
\[
H^2(X,\Lambda)\simeq H^1(\eta,\Lambda)
\simeq\Lambda(-1)\neq0.
\]
The last isomorphism follows from the valuation Kummer class of $t$; units are $\ell$-divisible by henselianity and algebraic closedness of $k$. Moreover $\cd_\ell(\eta)=1$ and the two charts have cohomological dimension zero, so the same Mayer--Vietoris sequence for arbitrary torsion sheaves gives $\cd_\ell(X)=2$.

The proposed value in (ii) is $\max(0,1)=1$, and that in the nonaffine clause of (iii) is $2\cdot1-1+0=1$. Thus both are false. This verifies every literal hypothesis, including finite type, integrality in (iii), and excellence of the base.
\end{proof}

\begin{theorem}[Proved from Gabber's affine theorem; A.7(i) and affine (iii)]
\label{A:relative-affine}
Let $Y=\Spec A$ be strictly local noetherian, let $f:X\to Y$ be of finite type, and let $\ell$ be invertible on $Y$. For $x\in X$ with image $y$, set
\[
\Delta(x)=\dim\overline{\{y\}}+\trdeg_{\kappa(y)}\kappa(x),
\qquad D=\sup_{x\in X}\Delta(x).
\]
If $Y$ is quasi-excellent and $X$ is affine, then
\[
\cd_\ell(X)\leq D\leq\max(\dim X_0,\dim X'+1).
\]
If in addition $X$ is integral and dominates $Y$, with generic relative dimension $d$, this gives $\cd_\ell(X)\leq\dim Y+d$.

For arbitrary $X$ and arbitrary noetherian strictly local $Y$, a valid replacement is $\cd_\ell(X)\leq2D$. In the quasi-excellent case this implies
\[
\cd_\ell(X)\leq2\max(\dim X_0,\dim X'+1).
\]
For integral dominant $X$, even without quasi-excellence, it implies
$\cd_\ell(X)\leq2\dim Y+2d$. The doubled trait attains this latter bound, so the missing unit in A.7(iii) cannot be recovered without an additional hypothesis.
\end{theorem}
\begin{proof}
We first establish the dimension bookkeeping. For a specialization $x\leadsto z$ in $X$, with images $y\leadsto w$, put $Z=\overline{\{x\}}$ with its reduced structure and $c=\dim\OO_{Z,z}$. The dimension inequality for finite type domain extensions, \citet[Tag 02IJ, Lemma 10.113.1]{A_Stacks}, gives
\[
c\leq\dim\OO_{\overline{\{y\}},w}
 +\trdeg_{\kappa(y)}\kappa(x)
 -\trdeg_{\kappa(w)}\kappa(z).
\]
Concatenating prime chains in the local domain $A/\mathfrak p_y$ gives
\[
\dim\OO_{\overline{\{y\}},w}
+\dim\overline{\{w\}}
\leq\dim\overline{\{y\}}.
\]
Thus $\Delta(x)-\Delta(z)\geq c$. In particular the supremum $D$ is the maximum of $\Delta$ at the finitely many generic points of $X$. If $X$ is integral dominant, this maximum is $\dim Y+d$.

When $Y$ is quasi-excellent it is excellent, since it is henselian. The dimension inequalities above are equalities on integral closures of points; consequently $y\mapsto\dim\overline{\{y\}}$ is a dimension function and its pullback is $\Delta$. Gabber's affine theorem, \citet[Exposé XV, Corollaire 1.2.2]{A_ILO2014}, therefore gives $\cd_\ell(X)\leq D$ for affine $X$. The extension from constructible coefficients is Lemma~\ref{A:basic}.

To bound $D$ in terms of the two fibres, consider a generic point $x$ of an irreducible component $Z$ of $X$. If its image is the closed point of $Y$, then $Z\subset X_0$ and $\Delta(x)=\dim Z\leq\dim X_0$. Otherwise put $R=A/\mathfrak p_y$, $e=\dim R>0$, and $r=\trdeg_{\Frac R}\kappa(x)$. Choose an affine open of $Z$ containing $x$. Relative Noether normalization, after inverting a nonzero $b\in R$, makes this affine open finite and dominant over $\A^r_{R[1/b]}$. We may multiply $b$ by a nonzero element of the maximal ideal to arrange that $D(b)$ misses the closed point.

The excellent local domain $R$ is catenary, and $\dim R[1/b]=e-1$. Here is the needed elementary justification. Complete $b$ to a system of parameters $b,b_2,\ldots,b_e$. A minimal prime $\mathfrak q$ over $(b_2,\ldots,b_e)$ has height at most $e-1$, so is not the maximal ideal and does not contain $b$. In $R/\mathfrak q$ the element $b$ cuts out the closed point, whence the principal ideal theorem gives $\dim R/\mathfrak q=1$. Catenarity in the local domain gives $\operatorname{ht}\mathfrak q=e-1$. This supplies the lower bound for $\dim R[1/b]$; the upper bound follows because its primes are not maximal in $R$. The finite normalization map now exhibits an open part of $Z\cap X'$ of dimension $e-1+r=\Delta(x)-1$. Therefore $\Delta(x)\leq\dim X'+1$, proving the asserted estimate for $D$.

For the bound $2D$ we do not assume excellence or affineness. The field $\kappa(y)$ is the fraction field of the strictly henselian local domain $A/\mathfrak p_y$, so Theorem~\ref{A:strict-field} gives $\cd_\ell\kappa(y)=\dim\overline{\{y\}}$. The finite-type field formula then gives $\cd_\ell\kappa(x)=\Delta(x)$. There is no orderable-field exception: if $\ell=2$, a square root of $-1$ already lies in the strictly henselian base. Apply \citet[Exposé X, Théorème 4.1]{A_SGA4} to the function $\phi(x)=2\Delta(x)$. Its first condition is the field bound just proved. For a strict specialization $x\leadsto z$, the fraction fields of the strict localization of $\overline{\{x\}}$ at $z$ have cohomological dimension at most $c$, by Theorem~\ref{A:strict-field} on their local-domain components. Since $c\geq1$, the preceding inequality gives
\[
\cd_\ell\Frac\bigl(\OO_{Z,z}^{\mathrm{sh}}\bigr)
\leq c<2c\leq\phi(x)-\phi(z).
\]
These are precisely the second conditions of Artin's criterion. It yields $\cd_\ell(X)\leq\sup\phi=2D$. All claimed variants follow.
\end{proof}

\begin{proposition}[Known; a positive replacement for proper morphisms]
\label{A:proper-fibre}
For a proper morphism $f:X\to Y$ with $Y$ strictly local noetherian and $\ell$ invertible on $Y$,
\[
\cd_\ell(X)\leq2\dim X_0.
\]
No quasi-excellence hypothesis is needed.
\end{proposition}
\begin{proof}
Proper base change and exactness of global sections on the strictly local base identify $H^q(X,F)$ with $H^q(X_0,F|_{X_0})$ for every torsion sheaf $F$. The closed fibre is proper over a separably closed field. Its $\ell$-cohomological dimension is at most twice its dimension by \citet[Exposé X, Corollaire 4.3]{A_SGA4}; coefficient approximation gives the assertion for all torsion sheaves. This replacement applies to actual proper morphisms and does not cover the doubled trait.
\end{proof}

\begin{lemma}[Proved; induction through one-dimensional localizations]
\label{A:allresidues-lower}
Let $(A,\mathfrak m,k)$ be a noetherian local domain of dimension $d$, with fraction field $K$. If $\ell=2$, suppose every residue field of $A$ is nonorderable. Then
\[
d_\ell(K)\geq d+d_\ell(k).
\]
\end{lemma}
\begin{proof}
We argue by induction on $d$. The case $d=0$ is equality. Choose a prime $\mathfrak p$ of height one with $\dim(A/\mathfrak p)=d-1$. Such a prime is the first nonzero prime in a chain of length $d$ from $(0)$ to $\mathfrak m$: if it had height greater than one, inserting a longer chain below it would contradict the definition of $d$. This argument needs no catenarity assumption.

The integral closure $B$ of the one-dimensional noetherian domain $A_{\mathfrak p}$ in $K$ is a Dedekind domain; its spectrum is surjective over $\Spec A_{\mathfrak p}$, and every residue extension is finite. This is the Krull--Akizuki consequence \citet[Tag 09IG, Lemma 10.120.18]{A_Stacks}. Choose a maximal ideal $\mathfrak q$ of $B$. Then $V=B_{\mathfrak q}$ is a discrete valuation ring with fraction field $K$ and residue field $k'$ finite over $\kappa(\mathfrak p)$.

Finite field extensions preserve $d_\ell$ in the present nonreal case. For finite $d_\ell$ this is \citet[\S2.2, final lemma on finite extensions]{A_GO2008}; in characteristic $\ell$, an infinite value is equivalent to infinite $\ell$-rank, which is unchanged by finite extensions. In other characteristics the finite-index subgroup theorem gives the same conclusion, with only the excluded orderable $2$-primary exception. Thus Lemma~\ref{A:valuation-lower} and the induction hypothesis for the local domain $A/\mathfrak p$ give
\[
d_\ell(K)\geq1+d_\ell(k')
=1+d_\ell(\kappa(\mathfrak p))
\geq1+(d-1)+d_\ell(k).
\]
Every residue field of $A/\mathfrak p$ is a residue field of $A$, so all induction hypotheses hold.
\end{proof}

\begin{lemma}[Proved; nonreal residues after henselization]
\label{A:hens-nonreal}
Let $(A,\mathfrak m,k)$ be a henselian local ring with nonorderable residue field. Then every residue field of $A$ is nonorderable.
\end{lemma}
\begin{proof}
If $\operatorname{char}k\neq2$, write $-1=\sum_{j=1}^s\bar a_j^2$ in $k$ and choose a nonzero $\bar a_1$. Lift $\bar a_2,\ldots,\bar a_s$ to $A$. The polynomial $T^2+1+\sum_{j=2}^s a_j^2$ has the simple root $\bar a_1$ modulo $\mathfrak m$. Hensel's lemma supplies a lift, expressing $-1$ as a sum of squares in $A$ itself.

If $\operatorname{char}k=2$, the polynomial $T^2+T+2$ has the simple root $0$ modulo $\mathfrak m$, and hence has a root $z\in A$. Direct calculation gives $(2z+1)^2=-7$, so
\[
-1=(2z+1)^2+1^2+1^2+1^2+1^2+1^2+1^2.
\]
In either case the same identity survives in every residue field. An orderable field cannot satisfy it.
\end{proof}

\begin{theorem}[Proved; A.8(ii), and a strengthened nonreal form of A.9]
\label{A:local-lower}
Let $(A,\mathfrak m,k)$ be any noetherian local domain with fraction field $K$, and let $\ell\neq\operatorname{char}K$. If $\ell=2$, assume $k$ nonorderable. Then
\[
\cd_\ell(K)\geq\dim A+d_\ell(k)
\geq\dim A+\cd_\ell(k).
\]
In particular A.8(ii) holds without the quasi-excellence assumption; henselianity is not necessary for this inequality.
\end{theorem}
\begin{proof}
Let $A^h$ be the henselization. It is noetherian, local, and faithfully flat over $A$, has the same residue field and dimension, and is a filtered colimit of étale $A$-algebras. Choose a minimal prime $\mathfrak q$ with $\dim(A^h/\mathfrak q)=\dim A$, and put $B=A^h/\mathfrak q$. The ring $B$ is a henselian local domain. Since $A$ is a domain and $A\to A^h$ is flat, contraction of $\mathfrak q$ is $(0)$; the resulting extension $L=\Frac B$ of $K$ is algebraic separable. Consequently $G_L$ is a closed subgroup of $G_K$, and $\cd_\ell(K)\geq\cd_\ell(L)$.

For $\ell\neq2$ there is no real-field restriction. For $\ell=2$, Lemma~\ref{A:hens-nonreal} shows that all residue fields of $B$ are nonorderable. Apply Lemma~\ref{A:allresidues-lower} to $B$. Since $\operatorname{char}L=\operatorname{char}K\neq\ell$, its left-hand side is $\cd_\ell(L)$. This proves the first inequality.

If $\operatorname{char}k\neq\ell$, the second inequality is equality by definition. In characteristic $\ell$, ordinary étale cohomological dimension is at most one. Kato dimension zero forces both a perfect field and the vanishing of the Artin--Schreier groups over every finite extension, which implies $\cd_\ell(k)=0$; otherwise $d_\ell(k)\geq1\geq\cd_\ell(k)$. This proves the second inequality. The finite-dimension hypothesis in A.8(ii) excludes orderability when $\ell=2$, so the theorem applies to precisely the case requested there.
\end{proof}

\begin{corollary}[Proved; the cancelled margin in A.4]
\label{A:p-rank-margin}
Let $A$ be a noetherian integral ring in which the prime $p$ is nonzero and not a unit, and let $K=\Frac A$. Then
\[
\cd_p(K)\geq1+\delta(V(p)).
\]
No quasi-excellence hypothesis is required.
\end{corollary}
\begin{proof}
For every maximal point $\mathfrak p$ of $V(p)$, the principal ideal theorem gives $\operatorname{ht}\mathfrak p=1$. The normalization-and-localization argument in Lemma~\ref{A:allresidues-lower} supplies a discrete valuation of $K$ whose residue field $k'$ is finite over $\kappa(\mathfrak p)$. That residue field has characteristic $p$, so is nonorderable when $p=2$. Lemma~\ref{A:valuation-lower} gives
\[
\cd_p(K)\geq1+d_p(k')\geq1+\delta(k')
=1+\delta(\kappa(\mathfrak p)).
\]
Here $\operatorname{char}K=0$, because a prime integer is nonzero and noninvertible in the domain. Taking the supremum over these maximal points proves the assertion.
\end{proof}

\begin{theorem}[Proved; the real-field ambiguity and regularity in A.9]
\label{A:realexception}
For a noetherian local domain $A$ with residue field $k$ and fraction field $K$, the inequality
\[
\cd_r(K)\geq\dim A+\cd_r(k)
\tag{A.9}\label{A:nine-ineq}
\]
holds outside the following two possible exceptions:
\begin{enumerate}[label=(\alph*)]
\item $r=\operatorname{char}K>0$;
\item $r=2$, $k$ is orderable, and $\vcd_2(k)=\cd_2(k(i))$ is finite.
\end{enumerate}
Thus ``real closed'' is not the correct interpretation of the ambiguous qualification on 31/023. If $A$ is regular, exception (b) is unnecessary. Neither the regularity condition in that last unrestricted real-residue assertion nor exception (a) can simply be discarded.
\end{theorem}
\begin{proof}
Theorem~\ref{A:local-lower} treats every nonreal residue field and every odd prime. It remains to consider $r=2$, $k$ orderable. In particular $\operatorname{char}k=0$, and $2$ is a unit in $A$.

First suppose $\vcd_2(k)=+\infty$. Form the finite étale $A$-algebra $A[i]=A[T]/(T^2+1)$. Its closed fibre is the field $k(i)$, so $A[i]$ is local; if it has several minimal primes, choose one whose quotient $B$ has dimension $\dim A$. The resulting local domain is finite over $A$ and has residue field $k(i)$. Theorem~\ref{A:local-lower}, applied to $B$, gives infinite $2$-cohomological dimension for its fraction field, which is a finite extension of $K$. Subgroup monotonicity forces $\cd_2(K)=+\infty$, and \eqref{A:nine-ineq} follows.

If $A$ is regular and $k$ is orderable, its completion is a complete equicharacteristic-zero regular local ring. A coefficient field identifies it with $k[[t_1,\ldots,t_d]]$. This power-series ring embeds in an iterated Laurent-series field over an ordered field, and its fraction field is therefore orderable. The injection $A\hookrightarrow\widehat A$ embeds $K$ in that orderable field. Hence $K$ is orderable and $\cd_2(K)=+\infty$, proving the assertion for regular $A$ as well.

To see why ``real closed'' is too restrictive even under the manuscript's catenary and universally Japanese assumptions, take
\[
A=\bigl(\Q[x,y]/(x^2+y^2)\bigr)_{(x,y)}.
\]
It is an excellent one-dimensional local domain, with residue field $\Q$ and fraction field $\Q(i)(x)$. Its residue field is orderable but not real closed; its virtual $2$-dimension is two, whereas its fraction field has $2$-dimension three. The right-hand side of \eqref{A:nine-ineq} is infinite, and the inequality fails. Replacing $\Q$ by $\R$ gives the manuscript's simpler example with $\cd_2(K)=1$. Thus regularity cannot be removed if the conclusion is to include all real residue fields.

Finally, let $A=\F_p[[t]]$ and $r=p$. This is a complete excellent regular henselian discrete valuation ring. Artin--Schreier theory gives $\cd_p(\F_p)=1$ and $\cd_p(\F_p((t)))=1$: the upper bounds are the characteristic-$p$ field theorem, and the nonzero classes are respectively a nontrivial finite constant-field extension and, for the latter field, its corresponding unramified extension. The claimed lower bound would be $1\geq2$, which is false. This proves that exception (a) is indispensable even for regular complete rings.
\end{proof}

\begin{remark}[Relation with the original Artin argument]
The invertible-prime, nonreal specialization argument is the one in \citet[Exposé X, Corollaires 2.4 and 2.5]{A_SGA4}. Theorem~\ref{A:local-lower} extends its valuation step across a residue characteristic equal to $r$ by using the opened Kato--Gabber--Orgogozo results. It should not be attributed to Artin's invertible-prime theorem alone. The stronger conclusion that only the closed residue field can be the real obstruction removes the catenary and universally Japanese restrictions in the manuscript's N.B.
\end{remark}

\subsection{A.8(i) and L4--L5: strictly local fields and affine opens}

\begin{theorem}[Known; A.8(i) without excellence]
\label{A:strict-field}
Let $B$ be a strictly henselian noetherian local domain of dimension $d$, with fraction field $K$, and let $\ell$ be invertible in $B$. Then
\[
\cd_\ell(K)=d.
\]
Moreover every nonempty proper open $U\subset\Spec B$ satisfies $\cd_\ell(U)\geq d$. If $B$ is quasi-excellent and $U$ is affine, equality holds.
\end{theorem}
\begin{proof}
The upper bound for $K$ is exactly \citet[Exposé XVIII$_B$, Corollaire 4.1.3]{A_ILO2014}, obtained from Théorème 4.1.1 with the local map $B\to B$ equal to the identity. That corollary has no excellence hypothesis. The reverse inequality is \citet[Exposé X, Corollaire 2.5]{A_SGA4}. The lower bound for a nonempty proper open is \citet[Exposé XVIII$_B$, Théorème 6.1]{A_ILO2014}. For a quasi-excellent base, the affine upper bound follows from \citet[Exposé XV, Corollaires 1.2.2 and 1.2.4]{A_ILO2014}, with the dimension function $b\mapsto\dim\overline{\{b\}}$; Lemma~\ref{A:basic} extends the bound from constructible finite coefficients to all torsion sheaves. Thus both bounds agree.
\end{proof}

\begin{proposition}[Proved; explicit one-divisor witness for L4]
\label{A:one-divisor}
Let $k$ be a separably closed field, let $A=k[[x,y]]$, let $\ell\neq\operatorname{char}k$, and put $U=\Spec A[1/x]$. Then $\cd_\ell(U)=2$. More explicitly, with
\[
D_1=V(y)\cap U,\qquad D_2=V(y-x)\cap U,
\qquad j:U\setminus(D_1\cup D_2)\hookrightarrow U,
\]
we have $H^2(U,j_!\Z/\ell)\simeq\Z/\ell(-1)$.
\end{proposition}
\begin{proof}
Both $D_1$ and $D_2$ are closed in $U$, are isomorphic to $\Spec k((x))$, and are disjoint: their intersection would force both $x=0$ and $x$ invertible. Let $\Lambda=\Z/\ell$ and let $i_a:D_a\hookrightarrow U$ be the closed immersions. There is an exact sequence of étale sheaves
\[
0\longrightarrow j_!\Lambda\longrightarrow\Lambda_U
\longrightarrow i_{1*}\Lambda\oplus i_{2*}\Lambda\longrightarrow0.
\]
The strictly local scheme $\Spec A$ and the strictly local regular divisor $V(x)$ have no positive étale cohomology with these constant coefficients. Absolute purity for the regular divisor $V(x)$ and the localization sequence therefore give
\[
H^1(U,\Lambda)\simeq\Lambda(-1),\qquad H^2(U,\Lambda)=0.
\]
Under this identification the generator is the Kummer class of $x$. The same Kummer class generates $H^1(D_a,\Lambda)\simeq\Lambda(-1)$ for each $a$, by the complete discrete valuation field formula. The restriction map in degree one is consequently the diagonal inclusion
\[
\Lambda(-1)\longrightarrow\Lambda(-1)\oplus\Lambda(-1),
\qquad c\longmapsto(c,c).
\]
The long exact sequence identifies $H^2(U,j_!\Lambda)$ with its cokernel, namely $\Lambda(-1)$. The upper bound two is Theorem~\ref{A:strict-field}. This supplies the missing sheaf in the one-boundary-component test; it does not claim that constant coefficients alone detect the cohomological dimension.
\end{proof}

